\documentclass[dvipdfmx,a4paper,titlepage,9pt]{jsarticle}

\input{{C:/Users/Lenovo/texlive/2023/Macros/macro}}

\usepackage{fancyhdr}
\pagestyle{fancy}
\fancyhead[L]{\bf ``Metric Space'' Class Notes}
\fancyhead[C]{}
\fancyhead[R]{HUFS(G) `23 M. sh}
\fancyfoot[L]{}
\fancyfoot[C]{\thepage}
\fancyfoot[R]{}

\title{``Metric Space'' Class Notes}
\begin{document}
\begin{multicols*}{2}
\Ltitle{0}{Backgrounds}

\ltitle{Basic Knowledge about Sets}\\
\hspace{12pt}{\small\Gray($X,\,Y$: sets)} \v{-6}
\begin{itemize}
\item[-] \stitle{Equality between sets}\\
         \;: $X=Y\;\defar\;X\sseq Y\;\wedge\;Y\sseq X$. \v6
\item[-] \stitle{Intersection}: $X\cap Y\equiv\{x\mid x\in X\;\wedge\;x\in Y\}$, \\
         \stitle{Union}: $X\cup Y\equiv\{x\mid x\in X\;\vee\;x\in Y\}$. \v6
\item[-] $X$ and $Y$ are \brk{disjoint}$\;\defar\;X\cap Y=\Emptyset$. \v{12}
\item[-] \stitle{Subtraction}: $X\setminus Y\equiv\{x\mid x\in X\;\wedge\;x\not\in Y\}$. \v6
\item[-] \stitle{Distribute Law}
   \begin{itemize}
   \setlength{\leftskip}{-10pt}
   \item[$\cdot$] $X\cup(Y\cap Z)=(X\cup Y)\cap(X\cup Z).$
   \item[$\cdot$] $X\cap(Y\cup Z)=(X\cap Y)\cup(X\cap Z).$ \v6
   \end{itemize}
\item[-] \stitle{De Morgan's Law}
   \begin{itemize}
   \setlength{\leftskip}{-10pt}
   \item[$\cdot$] $X\setminus(Y\cap Z)=(X\setminus Y)\cup(X\setminus Z).$
   \item[$\cdot$] $X\setminus(Y\cup Z)=(X\setminus Y)\cap(X\setminus Z).$ \v{12}
   \end{itemize}
\item[-] \stitle{Equality between Ordered pairs} \\
         \;: $(x,\;y)=(x',\;y')\;\defar\;x=x'\;\wedge\;y=y'$ \\
         \;\;\;{\Gray$(x,\,y,\,x',\,y'\in X)$} \v6
\item[-] \stitle{Cartesian Product} \\
         \;: $X\times Y\equiv\{(x,\;y)\mid x\in X\;\wedge\;y\in Y\}$ \\
         \;\;\;{\Gray$(\Emptyset\not\in\{X,\;Y\})$} \v6
\item[-] \stitle{Power set}
   \begin{itemize}
   \setlength{\leftskip}{-10pt}
   \item[$\cdot$] \brk{Power set of $X$}: $\msc{P}(X)\equiv\{U\mid U\subseteq X\}$.
   \item[$\cdot$] \brk{Intersection of $\msc{U\subseteq P}(X)$} \\
   \;: $\ds\bigcap\msc{U}\equiv\bigcap_{U\in\msc{U}}\!U\equiv\{x\mid \forall U\in\msc{U},\;x\in U\}$.
   \item[$\cdot$] \brk{Union of $\msc{U\subseteq P}(X)$} \\
   \;: $\ds\bigcup\msc{U}\equiv\bigcup_{U\in\msc{U}}\!U\equiv\{x\mid \exists U\in\msc{U},\;x\in U\}$.
   \end{itemize}
\end{itemize}

\vfill\null\columnbreak

\ltitle{Basic Knowledge about Functions}\\
\hspace{12pt}{\small\Gray($f$: function,\;$X,\,Y$: nonempty sets)} \v{-6}
\begin{itemize}[label=-]
\item\stitle{Definition of Function} \\
: For a set $F\subseteq X\times Y$, \\
\;\;if $\forall x\in X$,\;$\exists y\in Y$ s.t. $(x,\,y)\in F$, \\
\;\;then \brk{$f:=(X,\,Y,\,F)$} is \brk{a function from $X$ into $Y$}. \\

\item \stitle{Expression of Function} \\
: For a function $f=(X,\,Y,\,F)$ s.t. $(x,\,y)\in F$, we write\\
\;\;$f\!:\;\mt{X \longto Y, \v{-4} \\ x\lmapsto y,}$ or $y=f(x)$. \v5 \\
\;\;Also, for $U\subseteq X$, $f(U)$ is the set $\{f(x)\mid x\in U\}$. \\

\item\stitle{Sets Around Function} \\
 : For a function $f\!: X\longto Y,\;x\lmapsto y$, 
\begin{itemize}[label={$\cdot$}]
\item\brk{$X$} is the \brk{domain} of $f$.
\item\brk{$Y$} is the \brk{codomain} of $f$.
\item$\brk{f(X)}{\Gray\,(=\{f(x)\mid x\in X\})}$ is the \brk{range} of $f$.
\item$\brk{G(f)}{\Gray{}:=\{(x,\,f(x))\mid x\in X\}}$ is the \brk{graph} of $f$. \\
\end{itemize}

\item\stitle{Operation between Functions (1)} \\
{\small\Gray($f\!:X\longto\RR$, $g\!:Y\longto\RR$, $x\in(X\cap Y)$)}
\begin{itemize}[label={\Cdot}]
\setlength{\leftskip}{-10pt}
\item\brk{Addition/Subtraction} \\
: $\brk{f\pm g}\!:X\cap Y\longto\RR$ s.t. \\
\;\;$(f\pm g)(x)\equiv f(x)\pm g(x)$. \v4
\item\brk{Multiplication} \\
: $\brk{fg}\!:X\cap Y\longto\RR$ s.t. \\
\;\;$(fg)(x)\equiv f(x)\cdot g(x)$. \v4
\item\brk{Division} \v2\\
: $\brk{\ds\frac fg}\!: \{x\in X\cap Y\mid g(x)\neq0\}\longto\RR$ s.t. \\
\;\;$(f/g)(x)\equiv\ds\frac{f(x)}{g(x)}$. \\
\end{itemize}

\item\stitle{Types of Functions} \\
{\small\Gray($f\!:X\longto Y$, $x_1,\,x_2\in X$)}
\begin{itemize}[label={$\cdot$}]
\setlength{\leftskip}{-10pt}
\item$f$ is \brk{injective} if $f(x_1)=f(x_2)$ implies $x_1=x_2$.
\item$f$ is \brk{surjective} if $f(X)=Y$.
\item$f$ is \brk{bijective} if $f$ is injective and also surjective. \\
\end{itemize}

\item\stitle{Operation between Functions (2)} \\
{\small\Gray($f\!:X\longto Y$, $g\!:Z\longto W$, $(x,\,y)\in G(f)$)}
\begin{itemize}[label={\Cdot}]
\setlength{\leftskip}{-10pt}
\item\brk{Inverse Function} \\
: $\brk{f^{-1}}\!:Y\longto X$ s.t. $f^{-1}(y)=x$ exists \\
\;\;if and only if $f$ is bijective. \v4
\item\brk{Composed Function} \\
: $\brk{g\Circ f}\!:X\longto W$ s.t. $(g\Circ f)(x)=g(f(x))$ exists \\
\;\;if $f(X)\subseteq Z$.
\end{itemize}
\end{itemize}

\newpage

\Ltitle{1.1}{Metric Fcn. \& Open Ball}

\ltitle{Metric Functions}\v{-6}
\begin{itemize}[label=-]
\item\stitle{Definition} \\
: For a nonempty set $X$ and $\RR^+=[0,\,+\infty)$, a function\v{-4}
$$d\!:X\times X\longto\RR\v{-4}$$
\;\;is a \brk{metric on $X$} if it satisfies following 3 properties:
\begin{itemize}
\setlength{\leftskip}{15pt}
\item[(M1)] $d(x,\,y)=0$ iff $x=y$,
\item[(M2)] $d(x,\,y)=d(y,\,x)$,
\item[(M3)] $d(x,\,z)\lqs d(x,\,y)+d(y,\,z)$
\end{itemize}
\;\;for all $x,\,y,\,z\in X$.\v4
\item A \brk{metric space} is a pair $(X,\,d)$ or $\ul{X}$.\v4
\item We call $x\in X$ \brk{the point of $X$}.\v4
\item The \brk{distance of $x,\,y$} is the image of $d$, i.e., $d(x,\,y)$. \v4
\item\stitle{Theorems} \\
{\small\Gray[$(X,\,d)$: m.space, $x,\,y,\,z,\,w\in X$, $n\in\NN$, $1\lqs i\lqs n+1$]\v2}
\begin{itemize}[label={$\cdot$}]
\setlength{\leftskip}{-10pt}
\item$d(x_i,\,x_{n+1})\lqs\ds\sum_{j=i}^nd(x_j,\,x_{j+1})$.
\item$|d(x,\,z)-d(y,\,z)|\lqs d(x,\,y)$
\item$|d(x,\,y)-d(z,\,w)|\lqs d(x,\,z)+d(y,\,w)$
\end{itemize}
\end{itemize}

\;

\ltitle{Examples of Metric Functions}\v{-6}
\begin{itemize}[label=-]
\item\stitle{$\bs{d_E}$: Euclidean metric on $\RR^n$} $(n\in\NN)$ \\
: For $x:=(x_1,\,x_2,\,\ldots,x_n)$ and $y:=(y_1,\,y_2,\,\ldots,\,y_n),$\v{-4}
$$d_E\!:\RR^n\times\RR^n\longto\RR\quad\text{s.t.}\v{-6}$$
$$d_E(x,\,y)=\ds\sqrt{\sum_{i=1}^n(x_i-y_i)^2}\v{-4}\quad\text{for all}\;x,\,y\in\RR\v{-4}$$
\;\;is an \brk{Euclidean metric on $\RR^n$}. \\

\item\stitle{$\bs{d_D}$: Discrete metric}\h1{\scriptsize\Gray: 이산 거리 함수} \\
: For a nonempty set $X$, \v{-4}
$$d_D\!:X\times X\longto \RR\quad\text{s.t.}\v{-6}$$
$$d_D(x,\,y)=\begin{cases}0&\text{if}\;x=y,\\1&\text{if}\;x\neq y\end{cases}\quad\text{for all}\;x,\,y\in X\v{-4}$$
\;\;is a \brk{discrete metric on $X$}.\\

\item\stitle{Metric on $V$} \\
: For a vector space $V$ and the norm operator $\nrm{\cdot}$,\v{-4}
$$d\!:V\times V\longto\RR\quad\text{s.t.}\v{-6}$$
$$d(\mb{u,\,v})=\lVert\mb{u-v}\rVert\quad\text{for all}\;\mb{u,\,v}\in V\v{-4}$$
is a \brk{metric on $V$}. \v{-4}\\
$$\Gray
  \left(\nrm{\cdot}\!:\begin{array}{l}
                      \text{\ts{(\UTF{2170})}}\;\nrm{\mb v}=0\Lrarrow\mb{v=0} \v2 \\
                      \text{\ts{(\UTF{2171})}}\;\mb{\nrm{u+v}\lqs\nrm u+\nrm v} \v2 \\
                      \text{\ts{(\UTF{2172})}}\;\nrm{k\mb v}=|k|\cdot\nrm{\mb v}
                      \end{array}\right)$$

\vfill\null\columnbreak

\item\stitle{Supremum Metric on $B(X)$} \\
: For $X\neq\Emptyset$ and $B(X)\equiv\{f\mid f\!:X\longto\RR,\;f\!:\text{bdd.}\}$,\v{-4}
$$d\!:B(X)\times B(X)\longto\RR^+\quad\text{s.t.}\v{-6}$$
$$d(f,\,g)=\sup_{x\in X}{}|f(x)-g(x)|\quad\text{for all}\;f,\,g\in B(X)\v{-6}$$
is a \brk{supremum metric on $B(X)$}. \\

\item\stitle{$\bs{d_M}$: Maximum Metric on $\RR^n$} \\
: For $x:=(x_1,\,x_2,\,\ldots,\,x_n)$ and $y:=(y_1,\,y_2,\,\ldots,\,y_n)$,\v{-4}
$$d_M\!: \RR^n\times\RR^n\longto\RR^+\quad\text{s.t.}\v{-6}$$
$$d_M(x,\,y)=\max{}\{|x_i-y_i|:1\lqs i\lqs n\}\v{-4}$$
is a \brk{maximum metric on $\RR^n$}. \\

\item\stitle{$\bs{d_S}$: Sum Metric on $\RR^n$} \\
: For $x:=(x_1,\,x_2,\,\ldots,\,x_n)$ and $y:=(y_1,\,y_2,\,\ldots,\,y_n)$,\v{-4}
$$d_S\!: \RR^n\times\RR^n\longto\RR^+\quad\text{s.t.}\v{-6}$$
$$d_S(x,\,y)=\sum_{i=1}^n|x_i-y_i|\v{-4}$$
is a \brk{sum metric on $\RR^n$}. \\

\item\stitle{Inequality Between $d_E$, $d_M$, $d_S$} \\
: For $x,\,y\in\RR^n$, $d_M(x,\,y)\lqs d_E(x,\,y)\lqs d_S(x,\,y)$. \\

\end{itemize}

\ltitle{Open Ball} \v{-6}
\begin{itemize}[label=-]
\item\stitle{Definition} \\
: A set $\brk{S^d(x,\,\epsilon)}\equiv\{y\in X:d(x,\,y)<\epsilon\}$ is \\
\;\;an \brk{open ball with the center $x$ and radius $\epsilon$ relative to $d$}, \\
\;\;for a m.space $(X,\,d)$, given $\epsilon>0$, and a point $x\in X$. \\

\item\stitle{Examples}
\begin{itemize}[label={\Cdot}]
\setlength{\leftskip}{-10pt}
\item$S^{d_E}(0,\,1)=(-1,\,1)$ in $(\RR,\,d_E)$; \v2
\item$S^{d_E}((0,\,0),\,1)=\{(x,\,y)\in\RR^2:\sqrt{x^2+y^2}<1\}$; \v2
\item$S^{d_S}((0,\,0),\,1)=\{(x,\,y)\in\RR^2:|x|+|y|<1\}$; \v2
\item$S^{d_M}((0,\,0),\,1)=\{(x,\,y)\in\RR^2:|x|<1,\,|y|<1\}$. \\
\end{itemize}
\end{itemize}

\newpage

\Ltitle{1.2}{Important Points \& Subset} \v{-3}\\
\h{36}\tb{\Large\color[HTML]{538135}\ul{of Metric Spaces}} \v7

\ltitle{Inner Points/Sets}\v{-6}
\begin{itemize}[label=-]
\item\stitle{Inner Points} \\
: $x\in X$ is an \brk{inner point of a set $B\subseteq X$}, \\
\;\;if $\exists\epsilon>0$ s.t. $S(x,\,\epsilon)\subseteq B$. \\

\item\stitle{Inner Sets} \\
: For $B\in X$, $\inner B$ is an \brk{inner set of $A$}, \\
\;\;where $\brk{\inner{B}\equiv\{x\in X:x\;\text{is an inner point of $B$}\}}$.\v2\\
\;\;Thus, $x\in \inner{B}$ implies that $x$ is an inner point of $B$. \\
\end{itemize}

\ltitle{Open Set}\v{-6}
\begin{itemize}[label=-]
\item\stitle{Open Sets} \\
: For $B\subseteq X$, $B$ is an \brk{open set}, if \brk{$B=\inner{B}$}. \\
\;\;If $B$: open set, then $\forall x\in B$, $\exists\epsilon>0$ s.t. $S(x,\,\epsilon)\subseteq B$. \\

\item\stitle{Open Balls are Open Sets} \v2 \\
: $S(x,\,\epsilon)$ is an open set, for all $x\in X$ and given $\epsilon>0$.
\begin{pfbox}
\begin{proof}[]
Let $\delta=\epsilon-d(x,\,y)$ for $y\in S(x,\,\epsilon)$, and $z\in S(y,\,\delta)$.
Then, $d(y,\,z)<\delta$, hence
$$\begin{aligned}
  d(x,\,z)&\lqs d(x,\,y)+d(y,\,z) \\
          &<d(x,\,y)+\delta \\
          &<d(x,\,y)+\epsilon-d(x,\,y)=\epsilon,
  \end{aligned}$$
thus, $z\in S(x,\,\epsilon)$, which implies $S(y,\,\delta)\subseteq S(x,\,\epsilon)$.

Therefore, we have $\FA y\in S(x,\,\epsilon)$; $\exists\delta$ s.t. $S(y,\,\delta)\subseteq S(x,\,\epsilon)$, which explains $S(x,\,\epsilon)$ is an open set.
\end{proof}
\end{pfbox}

\;

\item\stitle{Theorems (1)} {\Gray($A,\,B\subseteq X$)}
\begin{enumerate}[label=(\alph*)]
\item $\inner{\Emptyset}=\Emptyset$, $\inner X=X$. \v2
\item $\inner A\subseteq A$. \v2
\item $B\subseteq A$ implies $\inner B\subseteq\inner A$. \v2
\item $\Inner{A\cap B}=\inner A\cap\inner B$. \v2
\item $\inner{\inner A}=\inner A$. \\
\end{enumerate}

\item\stitle{Theorems (2)}
\begin{enumerate}[label=(\alph*)]
\item $\Emptyset$, $X$ are open sets. \v2
\item If $A,\,B$ are open sets, then $A\cap B$ is so. \v2
\item If $\msc B$ is a set of open sets of $X$, \\
      $\ds\bigcup\msc{B}\equiv\{x\in B:\exists B\in\msc{B}\}$ is an open set.
\end{enumerate}

\vfill\null\columnbreak

\begin{pfbox}[]
\begin{proof}[(1d)]
Since $A\cap B\subseteq A$ \& $B$, we can conclude \v{-8}
$$\Inner{A\cap B}\subseteq \inner A\quad\tx{and}\quad\Inner{A\cap B}\subseteq \inner B\v{-8}$$
by (1c), which implies $\Inner{A\cap B}\subseteq\inner A\cap\inner B$. \v5

Let $x\in\inner A\cap\inner B$. Then, $\exists\epsilon_1,\;\epsilon_2>0$ s.t.\v{-4}
$$S(x,\,\epsilon_1)\subseteq A\quad\wedge\quad S(x,\,\epsilon_2)\subseteq B.\v{-4}$$

Put $\epsilon:=\min{}\{\epsilon_1,\,\epsilon_2\}$, and let $y\in S(x,\,\epsilon)$. Then $y\in S(x,\,\epsilon_1)$ and also $y\in S(x,\,\epsilon_2)$, thus $y\in S(x,\,\epsilon_1)\cap S(x,\epsilon_2)$, which is subset of $A\cap B$.

Hence, $y\in A\cap B$, then $S(x,\,\epsilon)\subseteq A\cap B$, thus $\inner A\cap\inner B\subseteq\Inner{A\cap B}$. \v5

Therefore, we have $\Inner{A\cap B}=\inner A\cap\inner B$.
\end{proof}\v{-5}
\dotfill
\begin{proof}[(1e)]
Let $x\in\inner A$. Then we have\v{-4}
$$\begin{aligned}
            & \tx{$\exists\epsilon>0$ s.t. $S(x,\,\epsilon)\subseteq A$;} \\
  \Longto{} & \Inner{S(x,\,\epsilon)}\subseteq\inner A\tx{;} \\
  \Longto{} & S(x,\,\epsilon)\subseteq\inner A\tx{;} \\
  \Longto{} & \tx{$\inner A$ is a open set;} \\
  \Longto{} & \inner{\inner A}=\inner A\tx{.}
  \end{aligned}\v{-20}$$
\end{proof}\v{-5}
\dotfill
\begin{proof}[(2c)]
By (1b), we have $\Inner{\ds\Union\msc{B}}\subseteq\ds\Union\msc B$.\hfill$\cdots\cdots\;(1)$

Let $x\in\ds\Union\msc B$. Then $x\in B$ for a open set $B\in\msc B$, thus $S(x,\,\epsilon)\subseteq\msc B$ for some $\epsilon>0$.

Then, $\exists\epsilon>0$ s.t. $S(x,\,\epsilon)\subset B\subset\ds\Union\msc B$.\hfill$\cdots\cdots\;(2)$

By $(1)$ and $(2)$, we have $\Inner{\ds\Union\msc{B}}=\ds\Union\msc B$, \\
which implies $\ds\Union\msc B$ is an open set.
\end{proof}
\end{pfbox}

\;

\item\stitle{Theorem (3)}\\
: For $B\subseteq X$, $\inner{B}$ is the largest open set contained by $B$.
\begin{Pfbox}
Let $A$ be a open set s.t. $A\subseteq B$. Then
\begin{itemize}
\item[(i)] $\inner B\subseteq B$ by (1b);
\item[(ii)] $\inner B$ is an open set by (1e);
\item[(iii)] $A\subseteq B\;\Longto\;\inner A\subseteq\inner B\;\Longto\;A\subseteq\inner B$;
\end{itemize}
thus we can conclude that $B$ is the largest open set contained in $B$.
\end{Pfbox}
\end{itemize}

\newpage

\ltitle{Distance of subsets}\v{-6}
\begin{itemize}[label=-]
\item\stitle{Definition} \\
: For $A,\,B\subseteq X$ where $\ul{X}$ is a metric space, \\
\;\;\brk{distance of $A$,\,$B$} is defined to a function\v{-4}
$$\dist\!: \P(X)\times\P(X)\longto\RR^+\cup\{\infty\},\v{-6}$$
such that\v{-4}
$$\h{-10}\dist(A,\,B)\equiv\begin{cases}\infty & \tx{if\;}A\tx{\;or\;}B=\Emptyset,\\
                    \inf{}\{d(a,\,b):a\in A,\,b\in B\} & \tx{if\;}A,\,B\neq\Emptyset.\end{cases}$$
                    \v{-6}
                            
\item\stitle{Properties} ($\ul{X}$: a metric space, $A,\,B,\,C\subseteq X$)\v2
\begin{itemize}
\setlength{\leftskip}{-5pt}
\item[(a)] $\dist(A,\,B)=\dist(B,\,A)$.\v2
\item[(b)] $A\subseteq C$ implies $\dist(A,\,C)\gqs\dist(B,\,C)$.\v2
\item[(c)] $A\cap B\neq\Emptyset$ implies $\dist(A,\,B)=0$.
\end{itemize}

\begin{pfbox}
\begin{proof}[(a)]
Let $p\!:\,A\tx{,}\,B\neq\Emptyset$,\;$q\!:\,a\in A,\,b\in B$. Then
$$\begin{aligned}
  \dist(A,\,B)&=\begin{cases}\infty &\tx{if\;}\neg p\tx{,}\\\inf{}\{d(a,\,b):q\} & \tx{if\;}p\end{cases} \\
              &=\begin{cases}\infty &\tx{if\;}\neg p\tx{,}\\\inf{}\{d(b,\,a):q\} & \tx{if\;}p\end{cases}\\
              &=\dist(B,\,A)\tx{.}
  \end{aligned}\v{-20}$$ 
\end{proof}\v{-5}
\dotfill
\begin{proof}[(b)]
If $B=\Emptyset$; $\dist(C,\,B)=\dist(A,\,B)=\infty$.

If $B\neq\Emptyset$ but $A=\Emptyset$; $\dist(A,\,C)=\infty\gqs\dist(B,\,C)$.

If $A,\,B\neq\Emptyset$, then\v{-4}
$$\begin{aligned}
  \dist(C,\,B)&   =\inf{}[D_C:=\{d(c,\,b):c\in C,\,b\in B\}] \\
              &\lqs\inf{}[D_A:=\{d(a,\,b):a\in A,\,b\in B\}] \\
              &   =\dist(A,\,B)\tx{,}
  \end{aligned}\v{-4}$$
since $D_C\subseteq D_A$.
\end{proof}\v{-5}
\dotfill
\begin{proof}[(c)]
By the assumption, there exists $c\in A\inter B$, thus\v{-4}
$$\begin{aligned}
  \dist(A,\,B)&   =\inf{}\{d(a,\,b):a\in A\tx{,}\,b\in B\} \\
              &\lqs d(c,\,c)=0\tx{,}
  \end{aligned}\v{-4}$$
which shows $\dist(A,\,B)=0$.
\end{proof}\v{-5}
\end{pfbox}
\end{itemize} \; \\

\ltitle{Neighbor Relation}\v{-6} \ti{\Gray(in short: nhb. \!rel.)}
\begin{itemize}[label=-]
\item\stitle{Definition} \\
: For $A$, $B\subseteq X$, \brk{$A$ and $B$ are in neighbor relation}, \\
\;\;if \brk{$\dist(A,\,B)=0$}. \v2\\
\;\;$A$ \& $B$ are in nhb. \!rel.,\; if and only if\\
\;\;\brk{$\forall\epsilon>0$; $\exists a\in A$, $b\in B$ s.t. $d(a,\,b)<\epsilon$}.

\vfill\null\columnbreak

\item\stitle{Properties} ($A$, $B$, $C\subseteq X$)
\begin{enumerate}[label=(\alph*)]
\item If $A$ \& $B$: nhb. \!rel.; $A$, $B\neq\Emptyset$. \v2
\item If $A$ \& $B$: nhb. \!rel.; $B$ \& $A$: nhb. \!rel. \v2
\item If $A$ \& $B$: nhb. \!rel. $\wedge$ $B\subseteq C$; \\
      $A$ \& $C$: nhb. \!rel. \v2
\item If $A$ \& $B\cup C$: nbh. \!rel; $A$ \& $B$ $\vee$ $A$ \& $C$: nhb. \!rel. \v2
\item If $A\cap B\neq\Emptyset$; $A$ \& $B$: nhb. \!rel.
\end{enumerate}
\begin{pfbox}
\begin{proof}[(d)]
By the assumption, $\dist(A,\,B\cup C)=0$\\
$\Longto\forall\epsilon>0$ $\exists x\in A$, $y\in (B\cup C)$ s.t. $d(x,\,y)<\epsilon$.

Since $y\in B_{\;\cdots\;(1)}$ or $y\in C_{\;\cdots\;(2)}$,\v2
\begin{itemize}
\setlength{\leftskip}{-5pt}
\item[$(1)$] $\dist(A,\,B)\lqs d(x,\,y)<\epsilon\Longto\dist(A,\,B)=0$,
\item[$(2)$] $\dist(A,\,C)\lqs d(x,\,y)<\epsilon\Longto\dist(A,\,C)=0$,\v2
\end{itemize}
Thus $A$ \& $B$: in nhb. rel. or $A$ \& $C$: in nhb. rel.
\end{proof}
\end{pfbox}\v9
\end{itemize}

\ltitle{Closedness}\v{-6}
\begin{itemize}[label=-]
\item\stitle{Adherent point} \ti{\Gray(in short: adh. \!pt.)}
\begin{itemize}[label={\Cdot}]
\setlength{\leftskip}{-5pt}
\item\stitle{Definition} \\
: For $A\subseteq X$ and $x\in X$, $x$ is an \brk{adherent point of $A$},\\
\;\;if \brk{$\dist(x,\,A)=0$}. \\

\item\stitle{Properties} ($A$, $B\subseteq X$, $x\in X$)
\begin{enumerate}[label=(\alph*)]
\setlength{\leftskip}{-5pt}
\item $x$: adh. \!pt. \!of $A$\quad iff\quad$\forall\epsilon>0$; $S(x,\,\epsilon)\cap A\neq\Emptyset$. \v2
\item If $x$: adh. \!pt. \!of $A$; $A\neq\Emptyset$. \v2
\item If $x$: adh. \!pt. \!of $A$ $\wedge$ $A\subseteq B$; \\
$x$: adh. \!pt. \!of $B$. \v2
\item If $x$: adh. \!pt. \!of $A\union B$; \\
$x$: adh. \!pt. \!of $A$ $\vee$ adh. \!pt. \!of $B$. \v2
\item If $x\in A$; $x$: adh. \!pt. \!of $A$. \v2
\item $x$: adh. \!pt. \!of $A$\quad iff\quad $\{x\}$ \& $A$: nhb. \!rel. \v5
\item[] \h{-20}\ti{\Gray(These are corollaries of the properties of nhb. rel.)}\\
\end{enumerate}
\end{itemize}

\item\stitle{Closure}
\begin{itemize}[label={\Cdot}]
\setlength{\leftskip}{-5pt}
\item\stitle{Definition} \\
: For $A\subseteq X$, $\ol{A}$ is the \brk{closure of $A$}, \\
\;\;where \brk{$\ol{A}\equiv\{x\in X:\tx{$x$ is an adh. \!pt. \!of $A$}\}$}. \v2\\
\;\;Thus, $x\in \ol{A}$ implies that $x$ is an adh. \!pt. \!of $A$. \\

\item\stitle{Properties} ($A$, $B\subseteq X$)
\begin{enumerate}[label=(\alph*)]
\setlength{\leftskip}{-5pt}
\item $\ol{\Emptyset}=\Emptyset$, $\ol X=X$. \v2
\item $A\subseteq \ol{A}$. \qquad {\Gray (cf. $\inner A\subseteq A$)} \v2
\item If $A\subseteq B$; $\ol A\subseteq \ol B$. \v2
\item $\ol{A\union B}=\ol A\union \ol B$. \v2
\item $\ol{\ol A}=\ol A$.
\end{enumerate}
\end{itemize}

\newpage

\begin{pfbox}
\begin{proof}[(b)]
Let $x\in A$. Then $\dist(x,\,A)=0$, thus $x\in\ol A$, \\
so we have $A\subseteq\ol A$.
\end{proof}\v{-5}
\dotfill
\begin{proof}[(d)]
Since $A\cup B\subseteq\ol A\cup\ol B$, $\ol{A\cup B}\subseteq\ol{\ol A\cup\ol B}=\ol A\cup\ol B$.

Also $A$, $B\subseteq\ol{A\cup B}$, thus $\ol A$, $\ol B\subseteq\ol{\ol{A\cup B}}=\ol{A\cup B}$.

Hence $\ol A\cup\ol B\subseteq\ol{A\cup B}$, we have $\ol A\cup\ol B=\ol{A\cup B}$.
\end{proof}\v{-5}
\dotfill
\begin{proof}[(e)]
By (b), $\ol A\subseteq\ol{\ol A}$ is obvious.\hfill $\cdots\cdots\;(1)$\v2

Let $x\in\ol{\ol A}$. Then we have $\dist(x,\,\ol A)=0$.\v2

Then, for $\epsilon>0$, $\exists y\in\ol A$ s.t. $d(x,\,y)<\epsilon/2$, \\
where $y$ is an adh. \!pt. \!of $A$. \v2

Thus $\dist(y,\,A)=0$, \;$\therefore\:\exists a\in A$ s.t. $d(a,\,y)<\epsilon/2$.

$$\begin{aligned}
  \Longto d(x,\,a)&\lqs d(x,\,y)+d(y,\,a) \\
                  &<    \epsilon/2+\epsilon/2=\epsilon\tx{,}
  \end{aligned}\v{-4}$$
which implies $\dist(x,\,A)=0$.\quad$\therefore\;x\in \ol A$.\hfill$\cdots\cdots\;(2)$\v2

By $(1)$ \& $(2)$, we have $\ol A=\ol{\ol A}$.
\end{proof}
\end{pfbox} \;

\item\stitle{Closed Set}
\begin{itemize}[label={\Cdot}]
\setlength{\leftskip}{-5pt}
\item\stitle{Definition} \\
: For $A\subseteq X$, $A$ is \brk{closed} / a \brk{closed set}, \\
\;\;if \brk{$\ol A=A$}.\\

\item\stitle{Properties} ($A$, $B$, $C\subseteq X$)
\begin{enumerate}[label=(\alph*)]
\item $\ol A$ is the smallest set which $A$ is a subset of.\v2
\item $\Emptyset$, $X$ are closed (and also open). \v2
\item If $A$, $B$ are closed; $A\cup B$ is so. {\Gray(cf. open: $\cap$)}\v2
\item If $\msc C$ is a set of closed sets of $X$; \\
      $\ds\bigcap\msc C\equiv\{x\in C:\forall C\in \msc C\}$ is also closed. \\
      {\Gray (cf. open: $\ds\bigcup\msc B\equiv\{x\in B:\exists B\in\msc B\}$)}
\end{enumerate}

\begin{pfbox}
\begin{proof}[(a)] Since
\begin{itemize}
\setlength{\leftskip}{-5pt}
\item[(i)] $A\subseteq \ol A$; (ii) $\ol{\ol A}=\ol A$ are obvious;\v2
\item[(ii)] For a closed set $B$ such that $A\subseteq B$,\\
            $\ol A\subseteq \ol B=B$;
\end{itemize}
$\ol A$ is the smallest closed set containing $A$.
\end{proof}\v{-7}

\dotfill\v2

\begin{proof}[(d)]
$\ds\Union\msc C\subseteq\ol{\ds\Union\msc C}$ is obvious. \hfill $(1)$
$$\begin{aligned}
           &\ds\Union\msc C\subseteq C\quad\tx{for all}\; C\in\msc C\tx{,} \\
  \Longto{}&\ol{\ds\Union\msc C}\subseteq \ol C\quad\tx{for all}\; C\in\msc C\tx{,}\\
  \Longto{}&\ol{\ds\Union\msc C}\subseteq\ds\Union_{C\in\msc C} C=\Union\msc C\tx{.}\quad\cdots\cdots\;(2)
  \end{aligned}$$
By $(1)$ \& $(2)$, we conclude $\ol{\ds\Union\msc C}=\ds\Union\msc C$.
\end{proof}
\end{pfbox}

\item\stitle{Relation of Inner \& Closure}
\begin{itemize}[label=-]
\setlength{\leftskip}{-10pt}
\item For $B\subseteq X$ and $x\in X$, \\
$\Longto$ \;\brk{$\inner{B}=X\setminus\ol{X\setminus B}$}. \v2\\
\qquad(\brk{$x$: inner pt. of $B$} \quad iff \quad \brk{$x\not\in\ol{X\setminus B}$}.)
\begin{Pfbox}
\v{-15}
$$\begin{aligned}
x\in\inner B &\ifff \exists\epsilon>0\tx{\;s.t.\;}S(x,\,\epsilon)\subseteq B\tx{;} \\
             &\ifff S(x,\;\epsilon)\cap (X\setminus B)=\Emptyset\tx{;}\\
             &\ifff x\not\in{X\setminus B}\tx{.}
\end{aligned}\v{-20}$$
\end{Pfbox}
\;\v{-5}
\item \brk{$B=\inner B$\quad iff \quad $X\setminus B=\ol{X\setminus B}$}. \v2\\
(\brk{$B$ is open\quad iff\quad $X\setminus B$ is closed.})
\begin{Pfbox}
\v{-15}
$$\begin{aligned}
B=\inner B &\ifff X\setminus B=X\setminus \inner B\tx{;} \\
           &\ifff X\setminus B=X\setminus(X\setminus\ol{X\setminus B})\tx{;} \\
           &\ifff X\setminus B=\ol{X\setminus B}\tx{.}
\end{aligned}\v{-20}$$
\end{Pfbox}
\;
\end{itemize}

\item\stitle{Theorems} ($A$, $B\subseteq X$)
\begin{itemize}[label=-]
\item If $\ol A\inter\ol B\neq\Emptyset$, then $A,\,B$ are in nhb. \!rel. \\
($\ifff$ \brk{$\ol A\inter\ol B\neq\Emptyset$ implies $\dist(A,\,B)=0$}.) \v5
\item \brk{$\dist(A,\,B)=\dist(\ol A,\,B)$}.\\
($\Longto$ \brk{$\dist(A,\,B)=\dist(\ol A,\,\ol B)$}.)
\end{itemize}
\end{itemize}
\end{itemize}

\;

\ltitle{Neighborhood} \ti{\Gray(in short: nbd.)}\v{-4}
\begin{itemize}[label=-]
\item\stitle{Definitions} ($A$, $U\subseteq X$)
\begin{itemize}[label=\Cdot]
\setlength{\leftskip}{-5pt}
\item\stitle{Neighborhood of a point} \\
: $U$ is a \brk{neighborhood of $x$}, if \brk{$x\in\inner U$}, \\
\;\;i.e., \brk{$\exists\epsilon>0$ s.t. $S(x,\,\epsilon)\subseteq U$}. \v4

\item\stitle{Neighborhood of a set} \\
: $U$ is a \brk{neighborhood of $A$}, if \brk{$a\in\inner U\;\,\forall a\in A$}, \\
\;\;i.e., \brk{for each $a\in A$; $\exists\epsilon>0$ s.t. $S(a,\,\epsilon)\subseteq U$}. \v4

\item\stitle{Uniform Neighborhood} \\
: $U$ is a \brk{uniform neighborhood of $A$},\\
\;\;if \brk{$\exists\epsilon>0$ s.t. $S(a,\,\epsilon)\subseteq U$ $\;\forall a\in A$}. \v6

\item[]\!\!\ti{\Gray(nbd. vs. uniform nbd. \v{-1}\\
\!\!\;\;: normal nbd. requires many $\epsilon$ for each $a\in A$,\v{-1}\\
\!\!\;\;\;\;but unifrom nbd. requires only single $\epsilon$ for all $a\in A$.)}
\end{itemize}

\newpage

\item\stitle{Theorems} ($A$, $U\subseteq X$, $x\in X$)\v2
\begin{itemize}[label=\Cdot]
\setlength{\leftskip}{-5pt}
\item $U$ is nbd. \!of $A$\quad$\Longto$\quad$A\subseteq U$. \v4

\item $a\in \inner A$\quad$\ifff$\quad$A$ is nbd. of $a$. \v4

\item $A$ is open\quad$\ifff$\quad$A$: nbd. of \;$\forall a\in A$.
\begin{Pfbox}
($\Longto$) Assume $A$ is open, and let $a\in A$. \\
then $a\in\inner A$, thus $A$ is nbd. of $A$.\hfill\Sq\v{-2}\\
\dotfill\v2\\
\;\,($\Longfrom$) Let $a\in A$, and assume $A$ is nbd. of $a$.\v2

Then, $a$ is an inner pt. of $A$, i.e., $a\in\inner A$, \\
which implies $A\subseteq\inner A$.\v2

And $\inner A\subseteq A$ is obvious, thus $\inner A=A$, \\
i.e., $A$ is open.
\end{Pfbox}\v4

\item $U$ is nbd. of $x$\quad$\ifff$\quad$x\not\in\ol{X\setminus U}$.
\begin{Pfbox}
($\Longto$) Assume $U$ is nbd. of $x$. \\
$\Longto$ $\exists\epsilon>0$ s.t. $S(x,\,\epsilon)\subseteq U$.\\
$\Longto$ $S(x,\,\epsilon)\cap(X\setminus U)=\Emptyset$.\v2\\
$\Longto$ $x\not\in\ol{X\setminus U}$.\hfill\Sq

\dotfill\v3

\;\,($\Longfrom$) Conversely, assume $x\not\in\ol{X\setminus U}$.\\
$\Longto$ $\exists\epsilon>0$ s.t. $S(x,\,\epsilon)\cap(X\setminus U)=\Emptyset$.\\
$\Longto$ $S(x,\,\epsilon)\subseteq U$. \\
$\Longto$ $U$: nbd. of $x$.
\end{Pfbox}\v9

\item \qquad\quad$U$ is uniform nbd. of $A$\\
$\ifff$\quad$A$ \& $X\setminus U$: not in nhb. rel.

\begin{pfbox}[($\Longto$)]
\begin{proof}[]
Assume $U$ is uniform nbd. of $A$.\\
$\Longto$ $\exists\epsilon>0$ s.t. $S(a,\,\epsilon)\subseteq U$ \;$\forall a\in A$.\v2\\
$\Longto$ $\ds\Union_{a\in A}S(a,\,\epsilon)\subseteq U$.\v2\\
$\Longto$ $\ds\Union_{a\in A}S(a,\,\epsilon)\cap(X\setminus U)=\Emptyset$.\hfill$\cdots\cdots\;(1)$\v2\\
$\Longto$ $\forall a\in A$, $\forall y\in X\setminus U$ s.t. $d(a,\,y)\gqs\epsilon$. \v4\\
\quad\;\;$\left(\begin{aligned}
&\tx{$\because$\;\;if not; $\exists a\in A$, $y\in X\setminus U$, $d(a,\,y)<\epsilon$,}\v{-2}\\
&\tx{which yields contradiction to (1).}
\end{aligned}\right)$ \v2\\
$\Longto$ $A$ \& $X\setminus U$ are not in nhb. rel.
\end{proof}
\end{pfbox}

\vfill\null\columnbreak

\begin{pfbox}[($\Longfrom$)]
\begin{proof}[]
Assume $A$ and $X\setminus U$ is not in nhb. rel.\\
$\Longto$ $\exists\epsilon>0$ s.t. $\dist(A,\,X\setminus U)=\epsilon$.\\
$\Longto$ $\forall a\in A$; $S(a,\,\epsilon)\subseteq U$.\v4\\
\quad\;\;$\left(\begin{aligned}
\tx{\;$\because$\;\;if not;}\quad\;\; &\tx{$\exists a\in A$ s.t. $S(a,\epsilon)\not\subseteq U$,}\\
\ifff{} & \tx{$S(a,\,\epsilon)\cap(X\setminus U)\neq\Emptyset$,}\\
\Longto{} & \tx{$\exists y\in X\setminus U$ s.t. $d(a,\,y)<\epsilon$,}\\
&\tx{contrd. to $\dist(A,\,X\setminus U)=\epsilon$.}
\end{aligned}\right)$\v4\\
$\Longto$ $U$ is uniform nbh. of $A$.
\end{proof}
\end{pfbox}\v9

\item If $U$ is (unif.) nbd. of $A$\; \&\; $B\subseteq A$\; \&\; $U\subseteq V$,\\
then $V$ is (unif.) nbd. of $B$.

\begin{Pfbox}
Let $b\in B$. Then, since $U$ is a nbd. of $A$, $\exists \epsilon>0$ s.t. $S(b,\,\epsilon)\subseteq U$.

Since $U\subseteq V$, $S(b,\,\epsilon)\subseteq U\subseteq V$, hence we have that $V$ is a nbd. of $B$.
\end{Pfbox}\v9

\item If $U$ \& $V$ are (unif.) nbd. of $A$, then is so $U\cap V$.

\begin{Pfbox}
(i) Assume $U$, $V$: nbd. of $A$, and let $a\in A$. \\
Then $\exists \epsilon_1$, $\epsilon_2>0$ s.t.\v{-4}
$$S(a,\,\epsilon_1)\subseteq U\quad\tx{and}\quad S(a,\,\epsilon_2)\subseteq V\tx{.}\v{-4}$$
Let $\epsilon=\min{}\{\epsilon_1,\,\epsilon_2\}$. Then we have\v{-4}
$$S(a,\,\epsilon)\subseteq U\quad\wedge S(a\,\epsilon)\subseteq V\tx{,}$$
which is equivalent to $S(a,\,\epsilon)\subseteq U\cap V$.\hfill\Sq

\dotfill\v3

(ii) Assume $U$, $V$: unif. nbd. of $A$, and let $a\in A$.\\
Then\v{-5}
$$\dist(A,\,X\setminus U)=\epsilon_1\tx{,}\quad\dist(A,\,X\setminus V)=\epsilon_2\v{-7}$$
for some $\epsilon_1$, $\epsilon_2>0$.\hfill$\cdots\cdots\;(1)$\v3

If $A$ \& $X\setminus(U\cap V)$ are in nhb. rel.,\\
$\Longto\;$ $(X\setminus U)\cup(X\setminus V)$ \& $A$ are in nhb. rel., \\
$\Longto\;$ $(X\setminus U)$ \& $A$: nhb. rel., or is $(X\setminus V)$ \& $A$,\\
but this contradicts to $(1)$.

Thus $A$ \& $X\setminus(U\cap V)$ is not in nhb. rel., hence we have\\
$U\cap V$ is unif. nbd. of $A$.
\end{Pfbox}\v7

\item If $U$ is unif. nbd. of $A$, then $U$ is unif. nbd. of $\ds\ol A$.
\begin{Pfbox}
By assumption, $\exists\epsilon>0$ s.t. $\dist(A,\,X\setminus U)=\epsilon$.

But, by Thm 1.2.28, $\dist(A,\,X\setminus U)=\dist(\ol A,\,X\setminus U)=\epsilon$,
hence $\ol A$ and $X\setminus U$ is not in nhb. rel.

Thus we have that $U$ is a unif. nbd. of $\ol A$. 
\end{Pfbox}
\end{itemize}
\end{itemize}

\newpage
\ltitle{Theorems about Neighborhoods}\v{-6}
\begin{itemize}[label=-]
\item If $|X|<+\infty$, then all subset $U$ of $X$ is clopen.
\begin{Pfbox}
Let $x\in X$, and let $y\in\ol{\{x\}}$.

Then $\dist(y,\,\{x\})=0$, thus $x=y$, $\ol{\{x\}}=\{x\}$.\v3

For $n=|X|$, $\ds\Union_{i=1}^{n-1}\{x_i\}=X\setminus\{x_1\}$ is closed,\\
hence $\ds X\setminus\Union_{i=1}^{n-1}\{x_i\}=\{x_1\}$ is open.

This implies $\{x_1\}$ is clopen, and so is $U\subseteq X$.
\end{Pfbox}\v7

\item For discrete metric space $(X,\,d_D)$,\v{-4}
$$\brk{\tx{$A$,\;$B$: in nhb. rel.}}\quad\ifff\quad \brk{A\cap B\neq\Emptyset}\tx{.}$$

\begin{Pfbox}
Since $A$ and $B$ are in nhb. rel., $\dist(A,\,B)=0$, thus for any $\epsilon>0$, $\exists a\in A$ and $b\in B$ s.t. $d_D(a,\,b)<\epsilon$.\v2

Since $\{d_D(a,\,b):a,\,b\in X\}=\{0,\,1\}$, and $\epsilon$ is arbitrary, we have $d_D(a,\,b)=0$, which implies $a=b$.\v2

Thus $\{a\}=\{b\}\subseteq A\cap B\neq\Emptyset$.
\end{Pfbox}\v7

\item For $A$,\;$B\subseteq\RR^n$, followings are equivalent:\v2
\begin{enumerate}[label=(\alph*)]
\item $A$,\;$B$ are in nhb. rel. with respect of $d_E$;\v1
\item $A$,\;$B$ are in nhb. rel. with respect of $d_M$;\v1
\item $A$,\;$B$ are in nhb. rel. with respect of $d_S$.\v1
\end{enumerate}
\begin{Pfbox}[{[(a)${}\Longto{}$(b)${}\Longto{}$(c)]}]
Note that $\FA x\tx{,}\,y\in \RR^n$,\v{-6}
$$d_M(x,\,y)\lqs d_E(x,\,y)\lqs d_S(x,\,y)\tx{,}\v{-3}$$
and let $x=(x_1,\,x_2,\,\ldots,\,x_n)$ and \\
$y=(y_1,\,y_2,\,\ldots,\,y_n)$ which belong to $\RR^n$.\v3

\dotfill\v7

\ti{[i: (a)${}\Longto{}$(b)]} \\
Assume $A$, $B$ are in nhb. rel. in $d_E$.\\
$\Longto\;\FA\epsilon>0$, $\exists(x,\,y)\in A\times B$ s.t. $d_E(x,\,y)<\epsilon$.\v2\\
$\therefore$ $d_M(x,\,y)\lqs d_E(x,\,y)<\epsilon$, \v2\\
thus $A$, $B$ are in nhb. rel. in $d_M$.\hfill\Sq\v3

\dotfill\v7

\ti{[ii: (b)${}\Longto{}$(c)]}\\
Assume $A$, $B$ are in nhb. rel. in $d_M$.\\
$\Longto$ $\FA\epsilon>0$, $\exists(x,\,y)\in A\times B$ s.t. $d_M(x,\,y)<\epsilon/n$.\v2

$\begin{aligned}
 &\Longto{}&&d_S(x,\,y)=\ds\sum_{i=1}^n|x_i-y_i| \\
 &&\lqs{}&n\cdot\max{}\{|x_i-y_i|:1\lqs i\lqs n\} \\
 &&<{}&n\cdot\epsilon/n=\epsilon\tx{.}
 \end{aligned}$\v2

Thus, $A$, $B$ are in nhb. rel. in $d_S$.
\end{Pfbox}

\vfill\null\columnbreak

\begin{Pfbox}[{[(c)${}\Longto{}$(a)]}]
Assume $A$, $B$ are in nhb. rel. in $d_S$.\\
$\Longto$ $\FA\epsilon>0$, $\exists(x,\,y)\in A\times B$ s.t. $d_S(x,\,y)<\epsilon$.\v2

Since $d_E(x,\,y)\lqs d_S(x,\,y)$ $\FA(x,\,y)\in A\times B$,\\
$d_E(x,\,y)\lqs d_S(x,\,y)<\epsilon$,\v2

thus $A$, $B$ is in nhb. rel. in $d_E$.
\end{Pfbox}
\end{itemize}

\;

\ltitle{Subspace}\v{-6}
\begin{itemize}[label=-]
\item\stitle{Definition}\\
: For a metric space $(X,\,d)$ and a subset $Y$ of $X$,\\
\;\;\brk{$(Y,\,d_{{\ds|}Y\times Y})$} is a \brk{subspace} of $(X,\,d)$ defined by $Y$.\\

\item\stitle{Properties} ($\ul X$: subspace of $\ul Y$, $A$, $B\subseteq X$, $x\in X$)\v1
\begin{itemize}[label=\Cdot]
\setlength{\leftskip}{-5pt}
\item $A$, $B$: nhb. rel. on $\ul X$\;${}\ifff{}$\;$A$, $B$: nhb. rel. on $\ul Y$.\v3
\item $x\in\ol A$ in $\ul X$\;${}\ifff{}$\;$x\in\ol A$ in $\ul Y$.\v3
\item $B$: (unif.) nbh. of $A$ in $\ul X$\\
$\ifff{}$\;$\exists U$: (unif.) nbh. of $A$ on $\ul Y$ s.t. $B=U\cap X$.\v3
\item $A$: open/closed in $\ul X$\\
$\ifff{}$\;$\exists O$: open/closed in $\ul Y$ s.t. $A=O\cap X$.
\end{itemize}
\end{itemize}

\; \\

\Ltitle{1.3}{Sequences in Metric Space}

\ltitle{Sequences}\v{-6}
\begin{itemize}[label=-]
\item\stitle{Definition}\\
: For a map $x\!:\,\NN\longto X$,\; $n\longmapsto x(n)=:x_n$,\\
\;\;the range of $x$, i.e., $(x_n):=\{x_n:n\in\NN\}$\\
\;\;is a sequence in $X$. \\

\item\stitle{Convergence of sequences}\v1
\begin{itemize}[label=\Cdot]
\setlength{\leftskip}{-5pt}
\item $(x_n)$ \brk{converges to $x$}, if \brk{for any $U$: nbd. of $x$,}\\
\brk{$\exists n\in\NN$ s.t. $m\gqs n$ implies $x_m\in U$}.\v4

\item We denoting that $(x_n)$ converges to $x$ by\v{-7}
$$\brk{(x_n)\longto x}\quad\tx{or}\quad \brk{x=\lim_{n\to\infty}x_n}\tx{.}$$

\item We call $x$ the \brk{limit (value)} or \brk{convergent value}, \\
if $(x_n)\longto x$.\\
\end{itemize}

\item\stitle{Accumulation point}\\
: $\forall U$: nbd. of $x$ and $\forall n\in\NN$, $\exists m\gqs n$ s.t. $x_m\in U$.\\

\item[]\ti{\Gray(Limit value vs. Accumulation point \v{-1}\\
\;\;: Limit: $|\{m:\exists U\tx{\;s.t.\;}x_m\neq U\}|\in\NN$, but\v{-1}\\
\;\;\;\;Acc. \!pt.: $|\{m:\exists U\tx{\;s.t.\;}x_m\neq U\}|\in\NN\cup\{+\infty\}$.)}
\end{itemize}

\newpage

\ltitle{Theorems}\v{-4}
\begin{itemize}[label=-]
\item \qquad\h{1.5}$(x_n)\longto x$ in $\ul X$; \\
$\ifff (d(x,\,x_n))\longto 0$ in $\ul \RR$; \\
$\ifff{}$$\FA\epsilon>0$; $\exists n\in\NN$ s.t. $m\gqs n$ implies $d(x,\,x_m)<\epsilon$.

\begin{Pfbox}[]
\ti{[(a)${}\Longto{}$(c)]} Let $\epsilon>0$. \\
Then $\exists n\in\NN$ s.t. $m\gqs n$ implies $x_m\in S(x,\,\epsilon)$.

Then we have $d(x,\,x_m)<\epsilon$ for all $m\gqs n$.\hfill\Sq

\dotfill\v2

\h4\ti{[(c)${}\Longto{}$(a)]} 

Let $U$ be a nbd. of $x$. Then $\exists\epsilon>0$ s.t. $S(x,\,\epsilon)\subseteq U$.

For $m$ in the assumption, $d(x,\,x_m)<\epsilon$, which is equivalent to $x_m\in S(x,\,\epsilon)$.

Hence, $m\gqs n$ implies $x_m\in U$, thus $(x_n)\longto x$.
\end{Pfbox} \v7

\item \qquad\h{1.5}$x$ is an acc. pt. of $(x_n)$ in $\ul X$; \\
$\ifff 0$ is an acc. pt. of $(d(x,\,x_n))$ in $\ul\RR$;\\
$\ifff{}$Given $\epsilon>0$; f.e. $n\in\NN$, $\exists m\gqs n$ s.t. $d(x,\,x_m)<\epsilon$.\\

\item If $(x_n)\longto x$,\\
$x$ is a \brk{unique acc. pt.} and also a \brk{unique limit value}.

\begin{Pfbox}
Let $\epsilon>0$. Then $\exists n\in\NN$ s.t. $d(x,\,x_m)<\epsilon/2$\\
for $m\gqs n$. Assume that $y\in\RR$ is an acc. pt. of $(x_n)$.\\
Then $\exists k\gqs n$ s.t. $d(y,\,x_k)<\epsilon/2$, thus\v{-4}
$$\begin{aligned}
  d(x,\,y)&\lqs d(x,\,x_k)+d(x_k,\,y) \\
  &<(\epsilon/2)\cdot 2=\epsilon\tx{.}
\end{aligned}\v{-5}$$
hence we have $d(x,\,y)<\epsilon$, \;$\therefore\;x=y$.
\end{Pfbox} \v9

\item If $(x_n)\longto x$ and $(y_n)\longto y$, \\
then $(d(x_n,\,y_n))\longto d(x,\,y)$.

\begin{Pfbox}
Let $\epsilon>0$. Then $\exists n\in\NN$ s.t. $m>n$ implies\v{-4}
$$\min{}\{d(x,\,x_m),\,d(y,\,y_m)\}<\epsilon/2\tx{.}\v{-4}$$
$\begin{aligned}
\Longto d(x_m,\,y_m)&\lqs d(x_m,\,x)+d(x,\,y)+d(y,\,y_m) \\
&<d(x,\,y)+\epsilon\tx{,}\v2\\
d(x,\,y)&\lqs d(x,\,x_m)+d(x_m,\,y_m)+d(y_m,\,y) \\
&<d(x_m,\,y_m)+\epsilon\tx{.}
\end{aligned}$

Hence we have\v{-4}
$$|d(x_m,\,y_m)-d(x,\,y)|<\epsilon\quad\tx{for\;}m\gqs n\tx{,}\v{-4}$$
thus $(d(x_n,\,y_n))\longto d(x,\,y)$.
\end{Pfbox} \v9

\vfill\null\columnbreak

\item For $a\gqs0$, if $d(x_n,\,y_n)\lqs a$ for each $n$, \\
and $(x_n)\longto x$ and $(y_n)\longto y$, then $d(x,\,y)\lqs a$.

\begin{Pfbox}
By the assumption, $(d(x_n,\,y_n))\longto d(x,\,y)$.

Thus for $\epsilon>0$, $\exists n\in\NN$ s.t. $m\gqs n$ implies\v{-4}
$$d(x,\,x_m)<\epsilon/2\tx{,}\quad d(y,\,y_m)<\epsilon_2\tx{.}\v{-4}$$
$\begin{aligned}
\Longto d(x,\,y)&\lqs d(x,\,x_m)+d(x_m,\,y_m)+d(y_m,\,y)\\
&<a+\epsilon\tx{.}
\end{aligned}\v2$

Since $\epsilon$ is arbitrary, we have $d(x,\,y)\lqs a$.
\end{Pfbox}\v9

\item \stitle{Subsequential Limits}\v1
\begin{itemize}[label=\Cdot]
\setlength{\leftskip}{-5pt}
\item Let $(x_{n_i})$ be a subsequence of $(x_n)$. Then,\\
\;(i) $(x_n)\longto x\quad\Longto\quad(x_{n_i})\longto x$. \\
\;(ii) $x$: acc. pt. of $(x_{n_i})\quad\Longto\quad$so is of $(x_n)$.\\

\item For $A\subseteq X$ and $x\in X$,\\
$\brk{x\in\ol A}\quad\tx{iff}\quad\brk{\exists (a_n)\in A\;\tx{s.t.}\;(a_n)\longto x}$.

\begin{Pfbox}
\ti{[\,$\Longto$]} Assume $x\in\ol A$, then $\dist(x,\,A)=0$.

Thus for each $n\in\NN$, $\exists a_n\in A$ s.t. \!$d(x,\,a_n)<1/n$.

Since $(1/n)\longto 0$, we have $(a_n)\longto x$.\hfill\Sq\v4

\h{-5}$\left(\begin{aligned}&\tx{$\because$ by the denceness of $\QQ$, $\exists n\in\NN$ s.t. $0<1/n<\epsilon$.}\\&\tx{Then, $m\gqs n$ implies $d(x,\,a_m)<1/m<1/n<\epsilon$.}\end{aligned}\right)$\v4

\dotfill\v3

\ti{[$\Longfrom$]} Let $\epsilon>0$, then $\exists n\in\NN$ s.t. $d(x,\,a_m)<\epsilon$ \\
for $m\gqs n$. Then, $\dist(x,\,A)<\epsilon$, \\
thus $\dist(x,\,A)=0$, \;$\therefore x\in\ol A$.
\end{Pfbox} \v9
\end{itemize}
\end{itemize}

\ltitle{Neighbor Relationship between Sequences}\v{-6}
\begin{itemize}[label=-]
\item\stitle{Definition}\\
: $(x_n)$, $(y_n)\in X$ are in \brk{nhb. rel.} if \brk{$(d(x_n,\,y_n))\longto 0$}.\\

\item\stitle{Equivalent Relation}
\begin{itemize}[label=\Cdot]
\setlength{\leftskip}{-10pt}
\item\stitle{Definition}\\
: For a set $X$, $R\subseteq X\times X$ is an \brk{equivalent relation} if
\begin{enumerate}[label=(\roman*)]
\setlength{\leftskip}{-15pt}
\item $(x,\,x)\in R$;\hfill\ti{\Gray[Reflexive Condition]}\v2
\item $(x,\,y)\in R\;\Longto\;(y,\,x)\in R$;\hfill\ti{\Gray[Symmetric Cd.]}\v2
\item $(x,\,y)$, $(y,\,z)\in R\;\Longto\;(x,\,z)\in R$.\!\!\hfill\ti{\Gray[Transitive Cd.]}\v{13}
\end{enumerate}

\item\stitle{Equivalentness of Nhb. Rel.}\\
: For $\msc S=\{(x_n):\tx{a seq. on $X$}\}$, let $R\subseteq \msc{S\times S}$ s.t.\v{-4}
$$\brk{R=\{(x_n,\,y_n):\tx{$(x_n)$, $(y_n)$ are in nhb. \!rel.}\}}\tx{.}\v{-4}$$
Then, \brk{$R$ is an equivalent relation in $\msc S$}.

\newpage

\begin{Pfbox}
(i) \ti{Reflexive Condition.}

Let $(x_n)$ be a seq. Then $(x_n)$ are in nhb. rel. with itself, \\
thus $(x_n,\,x_n)\in R$.\hfill\Sq

\dotfill\v3

(ii) \ti{Symmetric Condition.}

Let $(x_n,\,y_n)\in R$.\v1

$\Longto{}(x_n)$ \& $(y_n)$ are in nhb. rel.\v1

$\Longto{}(y_n)$ \& $(x_n)$ are in nhb. rel.\v1

[\,$\because\; d(x_n,\,y_n)=d(y_n,\,x_n)$.]\v1

$\Longto{}(y_n,\,x_n)\in R$.\hfill\Sq\v1

\dotfill\v3

(iii) \ti{Transitive Condition.}

Let $(x_n,\,y_n)$, $(y_n,\,z_n)\in R$.

Then $0\lqs d(x_n,\,z_n)\lqs d(x_n,\,y_n)+d(y_n,\,z_n)$, \\
hence $(x_n,\,z_n)\in R$.
\end{Pfbox}\v9
\end{itemize}
\item\stitle{Theorems}
\begin{itemize}[label=\Cdot]
\item For a seq. $(x_n)$ s.t. $\lim x_n=x$,\\
\brk{$(x_n)$ \& $(y_n)$ are in nhb. rel.}\quad iff\quad\brk{$(y_n)\longto x$}.\\

\item For $A$, $B\subseteq X$,\\
\brk{$A$ \& $B$: nhb. rel.}\quad iff\quad\begin{tabular}{l}\brk{$\exists (a_n)\subseteq A$, $(b_n)\subseteq B$}\\\brk{s.t. $(a_n)$ \& $(b_n)$: nhb. rel.}\end{tabular}
\begin{Pfbox}
\ti{[\,$\Longto$]} Since $\dist(A,\,B)=0$, for each $n\in\NN$,
$\exists a_n\in A$ and $b_n\in B$ s.t. $d(a_n,\,b_n)<1/n$.\v2

Hence $(a_n)$ and $(b_n)$ are in nhb. rel.\hfill\Sq

\dotfill\v3

\v4\ti{[$\Longfrom$]} Since $(a_n)$ \& $(b_n)$ are in nhb. rel, \\
for each $\epsilon>0$, $\exists n\in\NN$ s.t. $d(a_n,\,b_n)<\epsilon$.\v2

Then for all $\epsilon >0$, $\dist(A,\,B)\lqs d(a_n,\,b_n)<\epsilon$,\\
thus $\dist(A,\,B)=0$, and $A$, $B$ are in nhb. rel.
\end{Pfbox}\v9

\item For seq.s $(x_n)$, $(y_n)\subseteq X$,\v2\\
\brk{$(x_n)$ \& $(y_n)$ are in nhb. rel.}\quad if and only if\v2\\
\brk{for any $M\subset\NN$ s.t. $|M|=+\infty$,\\
$\{x_m:m\in M\}$ \& $\{y_m:m\in M\}$ are in nhb. rel.}
\end{itemize}
\end{itemize}

\vfill\null\columnbreak

\Ltitle{2.1}{Continuous \& Uniformly Cont. Ftn.}

\h{-9}(Let $\ul X=(X,\,d)$, $\ul{X'}=(X',\,d')$, $\ul{X''}=(X'',\,d'')$ be m.spaces.)\v6

\ltitle{Continuity}\v{-6}
\begin{itemize}[label=-]
\item\stitle{Definition} \\
: For a ftn. $f\!:\,X\longto X'$, $f$ is \brk{continuous} at $x$,\\
\;\;if \brk{$x\in\ol A$ implies $f(x)\in \ol{f(A)}$}.\\

\item\stitle{Judging Theorems} {\color{TGreen}(at a point)}\\
: \qquad\, $f\!:\,X\longto X'$ is cont. at $x$,\v2\\
\;\;$\ifff\;\forall U$: nbd. of $f(x)$; \brk{$f^{-1}(U)$ is nbd. of $x$},\v2\\
\;\;$\ifff\;$for $(x_n)$ in $\ul X$ s.t. $(x_n)\longto x$, \brk{$(f(x_n))\longto f(x)$},\v2\\
\;\;$\ifff\;$given $\epsilon>0$, $\exists\delta>0$ s.t.\\
\;\;\qquad \brk{$d(x,\,y)<\delta$ $\FA y\in X$ implies $d'\h{-1}(f(x),\,f(y))<\epsilon$}.\\
\;\;\qquad($\ifff\brk{f(S^d(x,\,\delta))\subseteq S^{d'}(f(x),\,\epsilon)}$.)

\begin{pfbox}
\rnc{\qedsymbol}{\Sq}
\begin{proof}[{[(a)${}\Longto{}$(d)]}]
Let $\epsilon>0$, and $x\in X$.

Put $B=X'\setminus S(f(x),\,\epsilon)$, and $A=f^{-1}(B)$.

Then, since\v{-6}
$$\begin{aligned}S(f(x),\,\epsilon)\cap f(A)&=S(f(x),\,\epsilon)\cap f(f^{-1}(B)) \\
&\subseteq S(f(x),\,\epsilon)\subseteq B=\Emptyset\tx{,}\end{aligned}\v{-6}$$
$S(f(x),\,\epsilon)\cap f(A)=\Emptyset$, thus $f(x)\notin\ol{f(A)}$.\v2

$\Longto x\neq\ol A$, since $f$ is cont. at $x$.\v2

$\Longto \exists\delta>0$ s.t. $S(x,\,\delta)\cap A=\Emptyset$.\v2

$\Longto S(x,\,\delta)\cap f^{-1}(X'\setminus S(f(x),\,\epsilon))=\Emptyset$.\v2

$\Longto S(x,\,\delta)\cap(X\setminus f^{-1}(S(f(x),\,\epsilon)))=\Emptyset$.\v2

$\Longto S(x,\,\delta)\subseteq f^{-1}(S(f(x),\,\epsilon))$.\v2

$\Longto f(S(x,\,\delta))\subseteq f(f^{-1}(S(f(x),\,\epsilon)))\subseteq S(f(x),\,\epsilon)$.\v2

$\Longto f(S(x,\,\delta))\subseteq S(f(x),\,\epsilon)$.\v{-4}
\end{proof}

\dotfill

\begin{proof}[{[(d)${}\Longto{}$(b)]}]
Let $U$ be a nbd. of $f(x)$ in $\ul{X'}$.\v2

$\Longto \exists\epsilon>0$ s.t. $S(f(x),\,\epsilon)\subseteq U$,\\
$\qquad\!\exists\delta>0$ s.t. $f(S(x,\,\delta))\subseteq S(f(x),\,\epsilon)$ by (d).\v2

$\Longto f(S(x,\,\delta))\subseteq U$.\v2

$\Longto S(x,\,\delta)\subseteq f^{-1}(U)$.\v2

$\Longto f^{-1}(U)$ is a nbd. of $x$.\v{-4}
\end{proof}

\dotfill

\begin{proof}[{[(b)${}\Longto{}$(c)]}]
Let $(x_n)$ be a seq. in $\ul{X}$ s.t. $(x_n)\longto x$.\v2

Let $U$ be a nbd. of $f(x)$, then $f^{-1}(U)$ is one of $x$.\v2

$\Longto \exists n\in\NN$ s.t. $m\gqs n$ implies $x_m\in f^{-1}(U)$.\v2

$\Longto f(x_m)\in f(f^{-1}(U))\subseteq U$ for all $m\gqs n$.\v2

$\Longto (f(x_n))\longto f(x).$\v{-4}
\end{proof}

\dotfill\rnc{\qedsymbol}{\Bsq}

\begin{proof}[{[(c)${}\Longto{}$(a)]}] Let $x\in\ol A$.\v2

$\Longto\exists (a_n)$ in $\ul X$ s.t. $(a_n)\longto x$ by Thm 1.3.10(*).\v2

$\Longto$ By (c), $(f(a_n))\longto f(x)$, \;$\therefore$ $f(x)\in\ol{f(A)}$ by (*).\\
\end{proof}
\end{pfbox}

\newpage

\item\stitle{Judging Theorems} {\color{TGreen}(on $\ul X$)}\\
: \qquad\, $f\!:\,\ul X\longto\ul X'$ is cont. ftn.,\v2\\
\;\;$\ifff\;\forall A$: open set on $\ul{X'}$, $f^{-1}(A)$: open on $\ul{X}$.\v2\\
\;\;$\ifff\;\forall B$: closed set on $\ul{X'}$, $f^{-1}(B)$: closed on $\ul{X}$.

\begin{pfbox}
\begin{proof}[{[(a)${}\Longto{}$(b)]}]\rnc{\qedsymbol}{\Sq}
For an open subset $A$ of $X'$, \\
$A$ is nbd. of all $f(x)\in A$ in $\ul{X'}$.\v3

By $f$ is cont., and Thm 2.1.2(b),\\
$f^{-1}(A)$ is nbd. of all $x\in f^{-1}(A)$ in $\ul X$.\v3

Hence we have $f^{-1}(A)$ is open on $\ul X$.
\end{proof}\v{-4}

\dotfill

\begin{proof}[{[(b)${}\Longto{}$(c)]}]\rnc{\qedsymbol}{\Sq}
For an closed subset $A$ of $X'$,\\
$B:=X'\setminus A$ is an open set on $\ul X$.\v3

By (b), $f^{-1}(B)$ is an open set, thus\\
$\ol{f^{-1}(A)}=\ol{X\setminus f^{-1}(B)}=X\setminus f^{-1}(B)=f^{-1}(A)$,\\
so we have $f^{-1}(A)$ is an open set on $\ul X$.
\end{proof}\v{-4}

\dotfill

\begin{proof}[{[(c)${}\Longto{}$(a)]}]
For $A\subseteq X$, $x\in X$, assume $f(x)\notin\ol{f(A)}$.\v3

Then $x\notin f^{-1}(\ol{f(A)})$, but since $A\subseteq f^{-1}(\ol{f(A)})$ and by (c), $f^{-1}(\ol{f(A)})$ is closed.\v3

Thus $\ol A\subseteq\ol{f^{-1}(\ol{f(A)})}=f^{-1}(\ol{f(A)})$, $x\notin\ol A$.\v3

By contrapositioning this,\\
we have $x\in \ol A\Longto f(x)\in\ol{f(A)}$, thus $f$ is cont.
\end{proof}
\end{pfbox}
\end{itemize}

\vfill\null\columnbreak

\ltitle{Uniform Continuity}\v{-6}
\begin{itemize}[label=-]
\item\stitle{Definition}\\
: A function $f\!:\,\ul X\longto\ul{X'}$ is \brk{uniformly continuous}, if\\
\;\;\brk{$A$, $B$: nhb. rel. in $\ul X$ $\Longto$ so is $f(A)$, $f(B)$ in $\ul{X'}$}.\v6

\item\stitle{Judging Theorem}\\
: \qquad\,$f\!:\,\ul X\longto\ul X'$ is unif. cont. ftn.\v2\\
\;\;$\ifff\;$ For $U$: unif. nbd. of $A$ in $\ul{X'}$,\\
\;\;\qquad\;$f^{-1}(U)$: unif. nbd. unif. nbd. of $f^{-1}(A)$ in $\ul X$.\v2\\
\;\;$\ifff\;$ For $(x_n)$, $(y_n)$: in nhb. rel. in $\ul X$,\\
\;\;\qquad\;$(f(x_n))$, $(f(y_n))$: in nhb. rel. in $\ul{X'}$.\v2\\
\;\;$\ifff\;$ Given $\epsilon>0$, $\exists\delta>0$ s.t.\\
\;\;\qquad\;$x$, $y\in X$ \& $d(x,\,y)<\delta$ implies $d'(f(x),\,f(y))<\epsilon$.

\begin{pfbox}
\rnc{\qedsymbol}{\Sq}
\begin{proof}[{[(a)${}\Longto{}\h{-1}$(c)]}]
For seqs. $(x_n)$, $(y_n)$in $\ul X$, \\
\h{-10}suppose $(f(x_n))$, $(f(y_n))$: not in nhb. rel. in $\ul{X'}$.\v3

Then $\exists M\subseteq\NN$ s.t. $|M|=+\infty$ and\\
\h{-10}$\{f(x_m):m\in M\}$\&$\{f(y_m):m\in M\}$: not in nhb. rel.\v3

By (a), $\{x_m:m\in M\}$\&$\{y_m:m\in M\}$: not nhb. rel,\\
\h{-10}thus $(x_n)$ \& $(y_n)$: not in nhb. rel.

By contrapositioning this, we can show (c) stands.\v{-3}
\end{proof}

\dotfill

\begin{proof}[{[(c)${}\Longto{}\h{-1}$(d)]}]
If $\neg$(d), for each $n\in\NN$, $\exists x_n$, $y_n\in X$ s.t.\v{-4}
$$d(x_n,\,y_n)<1/n\tx{\quad\&\quad}d'(f(x_n),\,f(y_n))\gqs\epsilon\v{-6}$$
\h{-10}for some $\epsilon>0$.\v3

Then $(f(x_n))$ \& $(f(y_n))$ are not in nhb. rel., \\
\h{-10}while $(x_n)$ \& $(y_n)$ are in nbh. rel., thus $\neg$(c) stands.

So we have $\neg$(d)${}\Longto{}\neg$(c), hence (c)${}\Longto{}$(d).\v{-3}
\end{proof}

\dotfill

\begin{proof}[{[(d)${}\Longto{}\h{-1}$(b)]}]
If $U$: unif. nbh. of $A$ in $\ul{X'}$,\\
\h{-10}then $\exists\epsilon>0$ s.t. $\dist(A,\,X'\setminus U)=\epsilon$.\hfill$\cdots\cdots\;(1)$\v2

By (d), for such $\epsilon$, $\exists\delta>0$ s.t.\v{-6}
$$d(x,\,y)<\delta\quad\tx{implies}\quad d'(f(x),\,f(y))<\epsilon\v{-6}$$
\h{-10}for all $x$, $y\in X$.\v2

Thus, $\forall x\in f^{-1}(A)$, $\forall y\in X\setminus f^{-1}(U)=f^{-1}(X'\setminus U)$,\\
\h{-10}since $d'(f(x),\,f(y))\not<\epsilon$ by $(1)$, $d(x,\,y)\gqs\delta$, \\
\h{-10}hence $\dist(f^{-1}(A),\,X\setminus f^{-1}(U))\gqs\delta>0$,\\
\h{-10}therefore we have $f^{-1}(A)$ is a unif. nbd. of $f^{-1}(U)$.\v{-3}
\end{proof}

\dotfill\rnc{\qedsymbol}{\Bsq}

\begin{proof}[{[(b)${}\Longto{}\h{-1}$(a)]}]
Let $A$, $B\subseteq X$,\\
\h{-10}and assume $f(A)$ \& $f(B)$ be not in nbh. rel.\v3

Then $U=X'\setminus f(B)$ is unif. nbh. of $f(A)$.\\
\h{-10}By (b), $f^{-1}(U)$ is unif. nbh. of $f^{-1}(f(A))$.\v3

Since $A\subseteq f^{-1}(f(A))$ and $f^{-1}(U)=X\setminus f^{-1}(f(B))$\\
\h{-10}$\subseteq X\setminus B$, $X\setminus B$ is unif. nbh. of $A$, and therefore,\\
\h{-10}$A$ and $B$ are not in nbh. rel.\v3

By contrapositioning this, we can show (a) stands.
\end{proof}
\end{pfbox}
\end{itemize}\v{-9}

\newpage

\ltitle{Various Continuous Function}\v{-6}
\begin{enumerate}[label=(\alph*)]
\item\stitle{Lipschitz Continuous Function}\\
: A \brk{Lipshitz continuous function} is defined by \\
\;\;a function $f\!:\ul X\longto\ul{X'}$ such that\v{-4}
$$\brk{d'(f(x),\,f(y))\lqs r\cdot d(x,\,y)}\;\;\tx{$\forall x$, $y\in X$, \;f.s. $r>0$\tx{.}}$$

\item\stitle{Continuity of Uniformly Continuous Function}\\
: \brk{Uniform coninuous functions are continuous}.\v9

\item\stitle{Composite Function}\\
: For (unif.) cont. ftn. $f\!:\ul X\longto\ul{X'}$ and $g\!:\ul{X'}\longto\ul{X''}$,\\
\;\;the \brk{composite function}\v{-6}
$$\brk{g\Circ f}\!:\ul X\longto\ul{X''}\tx,\quad x\longmapsto g(f(x))\v{-6}$$
\;\;is \brk{also (uniformly) continuous}.\v9

\item\stitle{Imbedding Function}\\
: For $\brk{X\subseteq Y}$, the \brk{imbedding function} is defined by a ftn.\v{-6}
$$\brk{f\!:\ul X\longto\ul Y},\quad\brk{x\longmapsto x}\tx.\v{-6}$$
\;\;An \brk{imbedding function is uniformly continuous}.\v9

\item\stitle{Restriction Function}\\
: For a subsp. $Y$ of $X$ and a (unif.) cont. ftn. $f\!:\ul X\longto \ul{X'}$,\\
\;\;the \brk{restriction function}\v{-6}
$$\brk{f_{\scalebox{0.8}{$|$}{\ds{}_Y}}}\!:\ul Y\longto\ul{X'}\tx,\quad y\longmapsto f(y)\v{-6}$$
\;\;is \brk{also (uniformly) continuous}.\v9

\item\stitle{Induced Function}\\
: For a (unif.) cont. ftn. $f\!:\ul X\longto\ul{X'}$\\
\;\;and a subspace $Y$ of $X'$ s.t. $f(X)\subseteq Y$,\v2\\
\;\;the \brk{induced function by $f$}\v{-6}
$$\brk{f^{\scalebox{0.8}{$|$}{\ds{}_Y}}}\!: \ul X\longto \ul Y\tx,\quad x\longmapsto f(x)\v{-6}$$
\;\;is \brk{also (uniformly) continuous}.
\end{enumerate}

\begin{pfbox}[\!\!: (a), (b)]
\begin{proof}[(a)] Let $f\!:\ul X\longto\ul{X'}$ be a Lipschitz cont. ftn.,\\
and $A$, $B$ be subsets of $X$ in nbh. rel., i.e., $\dist(A,\,B)=0$.

Then, for some $r>0$,\v{-4}
$$\begin{aligned}
\dist(f(A),\,f(B))&=\inf{}\{d'(f(x),\,f(y)):x\in A\tx,\,y\in B\}\\
&\gqs r\inf{}\{d(x,\,y):x\in A\tx,\,y\in B\}=0\tx,
\end{aligned}\v{-4}$$
$\therefore\;\dist(f(A),\,f(B))=0$, so $f(A)$ \& $f(B)$ are in nbh. rel.
\end{proof}\v{-4}

\dotfill

\begin{proof}[(b)] Let $A\subseteq X$, and $f\!: \ul X\longto\ul{X'}$ be unif. cont. function.

For $x\in\ol A$, $\dist(x,\,A)=0$, and by def. of unif. cont.,\\
$\dist(f(x),\,f(A))=0$, hence $f(x)\in\ol{f(A)}$.
\end{proof}\v{-4}
\end{pfbox}

\vfill\null\columnbreak

\begin{pfbox}[\!\!: (c), (d), (f)]
\begin{proof}[(c1. unif.)]
\rnc{\qedsymbol}{\Sq}Let $f\!:\ul X\longto\ul{X'}$, $g\!:\ul{X'}\longto\ul{X''}$ be uni. cnt. ftns.

For $A$, $B\subseteq X$ s.t. in nbh. rel. in $\ul X$,\\
then, $f(A)$ \& $f(B)$ are in nbh. rel. in $\ul{X'}$,\\
and also $g(f(A))$ \& $g(f(B))$ are in nbh. rel. in $\ul{X''}$.
\end{proof}\v{-4}

\dotfill

\begin{proof}[(c2. cont.)] Let $f\!:\ul X\longto\ul{X'}$, $g\!:\ul{X'}\longto\ul{X''}$ be cont. ftns,\\
and $A$ be an open set in $\ul{X''}$.

Since $g$ is cont., $g^{-1}(A)$ is open in $\ul{X'}$,\\
and because $f$ is cont., $f^{-1}(g^{-1}(A))$ is also open in $\ul X$.

Since $f^{-1}\Circ g^{-1}=(g\Circ f)^{-1}$,\\
$(g\Circ f)^{-1}(A)$ is open if $A$ is open, thus $g\Circ f$ is cont.
\end{proof}\v{-4}

\dotfill

\begin{proof}[(d)] For $x$, $y\in X$, $d'(f(x),\,f(y))=d'(x,\,y)=1\cdot d(x,\,y)$,\\
hence $f$ is Lipschitz cont., also unif. cont.
\end{proof}\v{-4}

\dotfill

\begin{proof}[(f: cont.)] Let $f\!:\ul X\longto\ul{X'}$ be a cont ftn.,\\
$\ul Y$ be a subspace of $\ul{X'}$ s.t. $f(X)\subseteq Y$ and $A\subseteq Y$ be open.

By Thm 1.2.40(d), $\exists$ an open (in $\ul{X'}$) set $O$ s.t. $A=O\cap Y$.

Since\v{-4}
$$\begin{aligned}
(f^{\scalebox{0.8}{$|$}{\ds{}_Y}})^{-1}(A)&=(f^{\scalebox{0.8}{$|$}{\ds{}_Y}})(O\cap Y)\\
&=(f^{\scalebox{0.8}{$|$}{\ds{}_Y}})^{-1}(O)\cap(f^{\scalebox{0.8}{$|$}{\ds{}_Y}})^{-1}(Y)\\
&=f^{-1}(O)\cap f^{-1}(Y)=f^{-1}(O)\cap X\tx,\end{aligned}\v{-4}$$
and $f^{-1}(O)$ and $X$ are open, hence $(f^{\scalebox{0.8}{$|$}{\ds{}_Y}})^{-1}(A)$ is open.\\
Hence we conclude $f^{\scalebox{0.8}{$|$}{\ds{}_Y}}\!:\ul X\longto\ul Y$ is cont.
\end{proof}\v{-4}

\dotfill

\begin{proof}[(f: unif.)] Similar way to (f: cont.) with following:\\
For a subspace $\ul Y$ of $\ul X$, let $A$, $B\subseteq Y$.

$B$ is unif. nbh. of $A$ on $\ul Y$ if and only if\\
$\exists$ a set $U\subseteq X$ s.t. unif. nbh. of $A$ on $\ul X$ and $U\cap Y=B$.
\end{proof}
\end{pfbox}

\newpage

\ltitle{Operational Function}\nc{\add}{\opn{add}}\nc{\mult}{\opn{mult}}\nc{\inv}{\opn{inv}}\v{-6}
\begin{enumerate}[label=(\alph*)]
\item$\brk{\add}\!:\,\ul{\RR^2}\longto\ul\RR$, $(x,\,y)\longmapsto x+y$: \brk{unif. cont.}\v4

\item$\brk{\mult}\!:\,\ul{\RR^2}\longto\ul\RR$, $(x,\,y)\longmapsto xy$: \brk{cont. (not unif.)}\v4

\item$\brk{\inv}\!:\,\ul{\RR\setminus\{0\}}\longto\ul\RR$, $x\longmapsto x^{-1}$: \brk{cont. (not unif.)}\v4

\item For a subsp. $\ul X$ of $\ul\RR$ s.t. $0\notin\ol X$,\\
\brk{$\inv_{\scalebox{0.8}{$|$}{\ds{}_X}}\!: \ul X\longto\ul\RR$} is \brk{unif. cont.}\v4

\item For a bounded subsp. $\ul X$ of $\ul{\RR^2}$,\\
\brk{$\mult_{\scalebox{0.8}{$|$}{\ds{}_X}}\!: \ul X\longto\ul\RR$} is \brk{unif. cont.}\v4
\end{enumerate}

\begin{pfbox}
\begin{proof}[(a)]
For $\epsilon>0$, let $\delta=\epsilon/2$, and assume\v{-4}
$$d_E((x,\,y),\,(x',\,y'))=\sqrt{(x-x')^2+(y-y')^2}<\delta\tx.\v{-4}$$
Then we observe both $|x-x'|$, $|y-y'|<\delta$.

We calculate\v{-4}
$$\begin{aligned}
d_E(x+y,\,x'+y')&=|x-x'+y-y'|\\
&\lqs|x-x'|+|y-y'|\\
&<2\delta=\epsilon\tx,
\end{aligned}\v{-4}$$
hence $\add$ is a unif. cont. ftn.
\end{proof}\v{-3}

\dotfill

\begin{proof}[(b)]
$\forall(x,\,y)\in\RR$ and $\forall\epsilon>0$, let $\delta=\min\left\{1,\,\dfrac\epsilon{|x|+|y|+1}\right\}$,\\
and assume $d_E((x,\,y),\,(x',\,y'))<\delta$.

Then we find $|x-x'|$, $|y-y'|<\delta$, and $|x'|<|x|+1$.

We derive\v{-4}
$$\begin{aligned}
|xy-x'y'|&\lqs|xy-x'y|+|x'y-x'y'|\\
&=|y||x-x'|+|x'||y-y'|\\
&<|y|\delta+(|x|+1)\delta\\
&=\delta(|x|+|y|+1)\lqs\epsilon\tx,
\end{aligned}\v{-4}$$
thus $\mult$ is a cont. ftn.\v4

But though $((n,\,0))$ and $((n,\,1/n))$ are in nhb. rel.,\\
$(\mult((n,\,0)))=(0)$ and $(\mult((n,\,1/n)))=(1)$ are not,\\
therefore $\mult$ is not a unif. cont. ftn.
\end{proof}\v{-3}

\dotfill

\begin{proof}[(c)]
For any $x\in\RR\setminus\{0\}$ and $\epsilon>0$, let $\delta=\ds\min\left\{\frac{|x|}2,\,\frac{\epsilon x^2}2\right\}$,\\
and assume $|x-y|<\delta$. Then $|y|\gqs|x|/2$.\\
(because if not, $|x|\lqs|x-y|+y<2(|x|/2)=|x|$; contrad.)

We conclude\v{-4}
$$\abs{\frac1x-\frac1y}=\frac{|y-x|}{|x||y|}<\frac\delta{|x|^2/2}=\epsilon\tx,\v{-4}$$
hence $\inv$ is a cont. ftn.\v4

But though $(1/n)$ and $(1/(2n))$ are in nbh. rel.,\\
$(\inv(1/n))=(n)$ and $(\inv(1/(2n)))=(2n)$ are not,\\
so $\inv$ is not a unif. cont. ftn.
\end{proof}
\end{pfbox}

\vfill\null\columnbreak

\begin{itemize}[label=-]
\item\stitle{Projection, Seq. Ftn.}
\begin{itemize}[label=\Cdot]
\setlength{\leftskip}{-10pt}
\item \stitle{Projections of $\RR^n$}\\
: For $i=1$, $2$, $\ldots$, $n$, the \brk{$i$-th projection of $\RR^n$}\v{-6}
$$\brk{p_i}\!:\,\ul{\RR^n}\longto\ul\RR\tx,\quad \brk{(x_1,\,x_2,\,\ldots,\,x_n)\longmapsto x_i}\v{-6}$$
is uniformly continuous.\v5

\item\stitle{Sequential Function}\\
: Each $f_i\!:\ul X\longto\ul\RR$ are (unif.) cont. if and only if\v{-4}
$$\h{-10}\brk{(f_1,\,\ldots,\,f_n)}\!:\ul X\longto\ul{\RR^n}\tx,\quad \brk{x\longmapsto(f_1(x),\,\ldots,\,f_n(x))}\v{-4}$$
is (unif.) cont.
\end{itemize}

\begin{pfbox}
\begin{proof}[{\color{TGreen}Projection.}]
For each $i\in[1,\,n]\cap\NN$,\v{-4}
$$\begin{aligned}
&|p_i(x_1,\,x_2,\,\ldots,\,x_n)-p_i(y_1,\,y_2,\,\ldots,\,y_n)|\\
={}&|x_i-y_i|\\
\lqs{}&d_E((x_1,\,x_2,\,\ldots,\,x_n),\,(y_1,\,y_2,\,\ldots,\,y_n))\tx,
\end{aligned}\v{-4}$$
hence $p_i$ is Lipschitz cont., thus it is unif. cont.
\end{proof}\v{-3}

\dotfill

\begin{proof}[{\color{TGreen}Seq. Ftn. ($\Longto$\!)}]\rnc{\qedsymbol}{\Sq}
For each $i\in[1,\,n]\cap\NN$,\\
$p_i\Circ (f_1,\,f_2,\,\ldots,\,f_n)=f_i$ is (unif.) cont.
\end{proof}\v{-3}

\dotfill

\begin{proof}[{\color{TGreen}Seq. Ftn. ($\Longfrom$)}]
Given $\epsilon>0$, for each $i\in[1,\,n]\cap\NN$, $\exists\delta_i>0$ s.t.\v{-4}
$$d(x,\,y)<\delta\quad\tx{implies}\quad |f_i(x)-f_i(y)|<\epsilon/\sqrt n\tx.\v{-4}$$
If we set $\delta=\min{}\{\delta_i:i=1\tx,\,2\tx,\,\ldots\tx,\,n\}$, then
$d(x,\,y)<\delta$ implies\v{-4}
$$\begin{aligned}
&d_E((f_1,\,f_2,\,\ldots,\,f_n)(x)-(f_1,\,f_2,\,\ldots,\,f_n)(y))\\
={}&\sqrt{\sum_{k=1}^n(f_k(x)-f_k(y))^2}<\sqrt{n\cdot\left(\frac\epsilon{\sqrt n}\right)^{\!2}}=\epsilon\tx,
\end{aligned}\v{-4}$$
thus $(f_1,\,f_2,\,\ldots,\,f_n)$ is unif. cont.
\end{proof}
\end{pfbox}\,

\item\stitle{Operations between Functions}
\begin{itemize}[label=\Cdot]
\setlength{\leftskip}{-10pt}
\item For (unif.) cont. ftn. $f\!: \ul X\longto\ul\RR$ and $g\!:\ul X\longto\ul\RR$,\\
$(f+g)\!: \ul X\longto\ul\RR$ is also (unif.) cont.\v4

\item For cont. ftn. $f\!:\ul X\longto\ul\RR$ and $g\!:\ul X\longto\ul\RR$,\\
$(f\cdot g)\!:\ul X\longto\ul\RR$ is also cont.\v4

\item For cont. ftn. $f\!:\ul X\longto\ul\RR$ s.t. $0\notin f(X)$,\\
$\dfrac1f\!:\ul X\longto\ul\RR$ is also cont.
\end{itemize}

\begin{Pfbox}
Observe that $f+g=\add\Circ(f,\,g)$,\v2\\
$f\cdot g=\mult\Circ(f,\,g)$ and $\dfrac1f=\inv\Circ f$,\v2\\
and each ftns are (unif.) cont. if $f$, $g$ is so.
\end{Pfbox}
\end{itemize}
\newpage

\Ltitle{2.2}{Equivalent Metric \& Isomorphism}

\ltitle{Definitions}\;\;{\Gray($f\!: \ul X\longto\ul{X'}$: bij. ftn.)}\v{-6}
\begin{itemize}[label=-]
\item\stitle{Isometry / Metric Isomorphism}\\
: $f$ satisfying \brk{$d(x,\,y)=d'(f(x),\,f(y))$}.\v9

\item\stitle{Uniform Isomorphism}\\
: $f$ s.t. both \brk{$f$ \& $f^{-1}$ are unif. cont.}\v9

\item\stitle{Topological Isomorphism}\\
: $f$ s.t. both \brk{$f$ \& $f^{-1}$ are cont.}\v9

\item\stitle{Isomorphic Space}\\
: $\ul X$ and $\ul{X'}$ are following,\\
\;\;if \brk{there exists}, between $\ul X$ and $\ul{X'}$,\v2\\
\;\;$\left\{\h{-4}\begin{tabular}{lll}m.isomorphism $f$&$\Longto$&\brk{metrically isomorphic};\\
u.isomorphism $f$&$\Longto$&\brk{uniformly isomorphic};\\
t.isomorphism $f$&$\Longto$&\brk{topologically isomorphic}.
\end{tabular}\right.$\v9

\item\stitle{Equivalent Metric}\\
: Metrics $d$ and $d'$ on $X$ are\\
\;\;\brk{\{uniformly, topologically\} equivalent}, if\v{-6}
$$\brk{\id_X\!:\;(X,\,d)\longto(X,\,d')}\v{-6}$$
\;\;is \brk{\{uniform, topological\} isomorphism}.\v9
\end{itemize}

\vfill\null\columnbreak

\ltitle{Theorems}\v{-6}
\begin{itemize}[label=-]
\item \brk{Isometry is uniform isomorphism},\\
and \brk{uniform isomorphism is topologically isomorphism}.\v9

\item For any $\ul X$, \brk{$\id_X\!:\ul X\longto\ul X$ is an isometry}.\v9

\item If \brk{$f\!:\ul X\longto\ul{X'}$ is \{m/u/t\} isomorphism},\\
then \brk{so is $f^{-1}\!:\ul{X'}\longto\ul X$}.\v9

\item If both $\brk f\!:\ul X\longto\ul{X'}$, $\brk g\!:\ul{X'}\longto\ul{X''}$ are \brk{\{m/u/t\} iso.},\\
then \brk{so is $g\Circ f\!:\ul X\longto\ul{X''}$}.
\end{itemize}

\begin{pfbox}\rnc{\qedsymbol}{\Sq}
\begin{proof}[\color{TGreen}(a: Isometry $\Longto$ unif. iso.)]\\
Assume $f\!:(X,\,d)\longto(X',\,d')$ is an isometry.\v4

Then $d'(f(x),\,f(y))=d(x,\,y)$ for any $x$, $y\in X$,\\
and also $d(f^{-1}(x'),\,f^{-1}(y'))=d'(x',\,y')$ for any $x'$, $y'\in X'$,\v4

Hence both $f$ and $f^{-1}$ are Lipschitz cont., also unif.
\end{proof}\v{-3}

\dotfill

\begin{proof}[\color{TGreen}(b)] $d(\id_X(x),\,\id_X(y))=d(\id_X^{-1}(x),\,\id_X^{-1}(y))=d(x,\,y)$.
\end{proof}\v{-3}

\dotfill

\begin{proof}[\color{TGreen}(c)] Trivial since $(f^{-1})^{-1}=f$.
\end{proof}\v{-3}

\dotfill

\begin{proof}[\color{TGreen}(d)] If $f$, $g$, $f^{-1}$, $g^{-1}$ are (unif.) cont.,\\
then so are $g\Circ f$ and $f^{-1}\Circ g^{-1}$, so trivial.
\end{proof}
\end{pfbox}\v9

\begin{itemize}[label=-]
\item \brk{\{m/u/t\} iso.} satisfies the \brk{equivalent relation}\\
on $\{\ul X\mid\tx{$\ul X$ is a metric space}\}$.\v9

\item Metrics $d$, $d'$ on $X$ are \brk{unif. isomorphic}, if \brk{$\exists\epsilon$, $\delta>0$} s.t.\v{-4}
\begin{center}
\brk{$x$, $y\in X$\quad implies\quad $\epsilon\cdot d(x,\,y)\lqs d'(x,\,y)\lqs\delta\cdot d(x,\,y)$}.
\end{center}
\end{itemize}

\begin{Pfbox}
Let $x$, $y\in X$. By the assumption, we have\v{-4}
$$d(x,\,y)\lqs\frac1\epsilon d'(x,\,y)\quad\tx{and}\quad d'(x,\,y)\lqs\delta\cdot d(x,\,y)\tx.\v{-4}$$
Then $\id_X\!: (X,\,d)\longto(X,\,d')$, $x\longmapsto x$ and\\
$\id_X^{-1}\!: (X,\,d')\longto(X,\,d)$, $x\longmapsto x$ are\\
both bij., and Lipschitz cont., so unif. cont.
\end{Pfbox}\v{15}

\newpage

\Ltitle{2.3}{\begin{tabular}{ll}The expansion of (unif.) cont. ftn.\v{-2}\\on ranges $\ul{[0,\,1]}$, $\ul{(0,\,1)}$, $\ul\RR$\end{tabular}}

\ltitle{(Unif.) Cont. Functions}\v{-6}
\begin{itemize}[label=-]
\item For any $A\subseteq X$ s.t. $A\neq\Emptyset$, a ftn.\v{-5}
$$\brk{f\!: \ul X\longto\ul\RR\tx,\quad x\longto\dist(x,\,A)}\v{-6}$$
is \brk{uniformly continuous}.

\begin{Pfbox}
For any $x$, $y\in X$ and $a\in A$,\\
$\dist(x,\,A)\lqs d(x,\,a)\lqs d(x,\,y)+d(y,\,a)$.\v2

Then $\dist(x,\,A)-d(x,\,y)\lqs d(y,\,a)$,\\
$\dist(x,\,A)-d(x,\,y)\lqs\dist(y,\,A)$,\\
$\dist(x,\,A)-\dist(y,\,A)=f(x)-f(y)\lqs d(x,\,y)$.

Similarly, $f(y)-f(x)\lqs d(y,\,x)$.\v2

Then $|f(y)-f(x)|=d_E(f(x),\,f(y))\lqs d(x,\,y)$,\\
hence $f$ is Lipschitx [\& unif.] cont.
\end{Pfbox}\v9

\item\stitle{Expansion of Indicator Function}\\
: For $A$, $B\subseteq X$ s.t. $A, B\neq\Emptyset$ and $A\cap B=\Emptyset$, if\v{-4}
$$f\!: \ul{A\cup B}\longto \ul{[0,\,1]}\tx,\quad x\longmapsto\brk{\begin{cases}1&\tx{if $x\in A$,}\\0&\tx{if $x\in B$}\end{cases}}\v{-4}$$
\;\;has a \brk{(unif.) cont. expansive ftn.}\v{-4}
$$g\!:\ul X\longto\ul{[0,\,1]}\tx,\quad x\longmapsto f(x)\tx,\v{-4}$$
\;\;then \brk{$\ol A\cap \ol B=\Emptyset$}. (i.e., \brk{$A$, $B$: not in nbh. rel.})

\begin{Pfbox}
If not, $\exists x\in\ol A\cap\ol B$,\\
thus $\exists(a_n)$, $(b_n)$ in $A$, $B$, resp. s.t. $a_n$, $b_n\longto x$.

But $g(a_n)=f(a_n)=0$: $f(x)=0$,\\
$g(b_n)=f(b_n)=1$: $f(x)=1$;\\
contrary to $f$. Thus $\ol A\cap\ol B=\Emptyset$.
\end{Pfbox}\v9

\item\stitle{Urysohn Lemma}\\
: For $A$, $B\subseteq X$ s.t. $A$, $B\neq\Emptyset$ and $\ol A\cap\ol B=\Emptyset$,\v{-4}
$$f\!:\ul X\longto\ul{[0,\,1]}\tx,\quad \brk{x\longmapsto\frac{\dist(x,\,A)}{\dist(x,\,A)+\dist(x,\,B)}}\v{-4}$$
\;\;is a \brk{well-defined function}, and satisfies following:\\
\quad\Cdot\;\brk{$f$ is continuous};\\
\quad\Cdot\;\brk{$f$ is unif. cont.} if and only if \brk{$A$, $B$: not in nbh. rel.}\v9
\end{itemize}

\vfill\null\columnbreak

\ltitle{Uniform Convergence}\v{-6}
\begin{itemize}[label=-]
\item\stitle{Definition}\\
: For a seq. $(f_n)$ of (unif.) cont. ftn. $f_n\!:\ul X\longto\ul{[0,\,1]}$\\
\;\;and a ftn. $f\!:\ul X\longto\ul{[0,\,1]}$,\\
\;\;\brk{$(f_n)$ uniformly converges to $f$} if $\forall\epsilon>0$, $\exists m\in\NN$ s.t.\v{-4}
\begin{center}\brk{$n\gqs m$ and $x\in X$\quad imply\quad $|f_n(x)-f(x)|<\epsilon$}.
\v5\end{center}

\item\stitle{Preservation of (unif.) cont.}\\
: For a seq. $(f_n)$ of \brk{(unif.) cont. ftns} $f_n\!:\ul X\longto\ul{[0,\,1]}$,\\
\;\;if \brk{$(f_n)$ unif. converges} to $f\!:\ul X\longto\ul{[0,\,1]}$,\\
\;\;then \brk{$f$ is also (unif.) cont.}

\begin{Pfbox}[(unif.)]
Let $f_n$ be unif. cont. ftns. s.t. unif. converges to $f$.

Then given $\epsilon>0$, $\exists n\in\NN$ s.t.\v{-4}
$$x\in X\quad\tx{implies}\quad |f_n(x)-f(x)|<\epsilon/3\tx.\v{-4}$$
Since $f_n$ is unif. cont., $\exists\delta>0$ s.t.\v{-4}
$$d(x,\,y)<\delta\;\tx\&\; x\tx,\,y\in X\quad\tx{imply}\quad |f_n(x)-f_n(y)|<\epsilon/3\tx.\v{-4}$$
If $d(x,\,y)<\delta$, then we have\v{-4}
$$\begin{aligned}
&|f(x)-f(y)|\\
\lqs{}&|f(x)-f_n(x)|+|f_n(x)-f_n(y)|+|f_n(y)-f(y)|\\
<{}&3\cdot \epsilon/3=\epsilon\tx,
\end{aligned}\v{-4}$$
hence $f$ is unif. cont.
\end{Pfbox}
\end{itemize}\v9

\ltitle{Expansion Theorem} ($\ul A$: closed subspace of $\ul X$)\v{-6}
\begin{itemize}[label=-]
\item Any (unif.) cont. ftn. $f\!:\ul A\longto\brk{\ul{[0,\,1]}}$\\
\brk{can be expanded} to a \brk{(unif.) cont.} ftn. $g\!:\brk{\ul X}\longto\ul{[0,\,1]}$.\v9

\item Any (unif.) cont. ftn. $f\!:\ul A\longto\brk{\ul{(0,\,1)}}$\\
\brk{can be expanded} to a \brk{(unif.) cont.} ftn. $g\!:\brk{\ul X}\longto\ul{(0,\,1)}$.\v9

\item Any (unif.) cont. ftn. $f\!:\ul A\longto\brk{\ul{\RR}}$\\
\brk{can be expanded} to a \brk{cont. ftn.} $g\!:\brk{\ul X}\longto\ul{\RR}$,\\
but \brk{cannot be to a unif. cont. ftn.}\v9

\item \brk{$\ul{(0,\,1)}$} and \brk{$\ul\RR$} are \brk{not uniformly isomorphic}.
\end{itemize}

\newpage

\Ltitle{3.1}{Cauchy Sq. \& Compl. of Metric Sp.}

\ltitle{Cauchy Sequence}\v{-6}
\begin{itemize}[label=-]
\item\stitle{Definition}\\
: $(x_n)$ is a \brk{Cauchy sequence}, if given $\epsilon>0$, $\exists n\in\NN$ s.t.\v{-4}
$$\brk{m\tx,\;m'\gqs n\quad\tx{implies}\quad d(x_m,\,x_{m'})<\epsilon}\tx.$$

\item\stitle{Boundedness of Sets}\\
: For $A\subseteq X$, \brk{$A$ is bounded set} if\v{-4}
$$\brk{\diam A}:=\begin{cases} 0 & \tx{if $A=\Emptyset$,}\\ \brk{\sup{}\{d(a,\,b):a\tx,\;b\in A\}} & \tx{if $A\neq\Emptyset$}\end{cases}\brk{<\infty}\tx.\v{-4}$$
\;\;{\Gray(``$\diam$'' stands for ``\ul{diam}eter.'')}\v9

\item\stitle{Judging Theorems}
: Followings are equivalent:
\begin{itemize}[label=\Cdot]
\setlength{\leftskip}{-10pt}
\item $(x_n)$ is a \brk{Cauchy sequence}.
\item \brk{$\ds\lim_{n\to\infty}\diam{}\{x_m:m\gqs n\}=0$}.
\item \brk{Any two subsequences} of $(x_n)$ are \brk{in nbh. rel.}\v9
\end{itemize}

\item\stitle{Theorems}
\begin{itemize}[label=\Cdot]
\setlength{\leftskip}{-10pt}
\item For a \brk{unif. cont. ftn.} $\brk f\!:\ul X\longto\ul Y$,\\
if \brk{$(x_n)$ is a Cauchy seq.} on $X$, then \brk{so is $(f(x_n))$}.\v4

\item If \brk{$(x_n)$ is a Cauchy sequence}, then \brk{$\{x_n:n\in\NN\}$ is bdd}.\v4

\item For seq.s $(x_n)$, $(y_n)$ in \brk{nbh. rel.},\\
\brk{$(x_n)$ is a Cauchy seq.}\quad$\ifff$\quad\brk{$(y_n)$ is a Cauchy seq.}\v4

\item \brk{$(x_n)\longto x$}\quad$\ifff$\quad\brk{$(x_n)$ is Cauchy seq. with acc. pt. $x$.}
\end{itemize}
\end{itemize}

\begin{pfbox}
\begin{proof}[\color{TGreen}(a)]
Let $(x_{n_i})$, $(x_{n_j})$ be subseq.s of $(x_n)$.

Then $(x_{n_i})$, $(x_{n_j})$ are in nbh. rel. Since $f$ is unif. cont.,\\
$(f(x_{n_i}))$, $(f(x_{n_j}))$ are also in nbh. rel.
\end{proof}\v{-5}

\dotfill\v{-2}

\begin{proof}[\color{TGreen}(b)]
By Thm 3.1.3(b), $\exists n\in\NN$ s.t. $\diam{}\{x_m:m\gqs n\}<1$.

Put $r=\max{}\{d(a_i,\,a_n):i=1\tx,\,2\tx,\,\ldots\tx,\,n\}$.

Then $\diam{}\{x_n:n\in\NN\}<2(r+1)\;(<+\infty)$.
\end{proof}\v{-5}

\dotfill\v{-2}

\begin{proof}[\color{TGreen}(c)]
Given $\epsilon>0$, $\exists n_0\in\NN$ s.t. $\diam{}\{x_m:m\gqs n_0\}<\epsilon/3$,\\
$\exists n_1\in\NN$ s.t. $n>n_1$ implies $d(x_n,\,y_n)<\epsilon/3$.

Let $m$, $n\gqs\max{}\{n_0,\,n_1\}$. Then we have\v{-4}
$$\begin{aligned}
d(y_n,\,y_m)&\lqs d(y_n,\,x_n)+d(x_n,\,x_m)+d(x_m,\,y_m)\\
&<3\cdot\epsilon/3=\epsilon\tx.
\end{aligned}\v{-20}$$
\end{proof}\v{-5}

\dotfill\v{-2}

\begin{proof}[\color{TGreen}(d\,$\To$)]\rnc{\qedsymbol}{\Sq}$\forall\epsilon>0$, $\exists n\in\NN$ s.t. $d(x_m,\,x)<\epsilon/2$ for $m\gqs n$.

For $m$, $m'\gqs n$, $d(x_m,\,x_{m'})\lqs d(x_m,\,x)+d(x,\,x_{m'})<\epsilon$.
\end{proof}\v{-5}

\dotfill\v{-2}

\begin{proof}[\color{TGreen}(d\,$\From$)] $\forall\epsilon>0$, $\exists n_0\in\NN$ s.t. $\diam{}\{x_m:m\gqs n\}<\epsilon/2$.

Since $x$: acc. pt., $\exists n_1\gqs n_0$ s.t. $d(x_{n_1},\,x)<\epsilon/2$.

Then, for $n\gqs n_0$, $d(x,\,x_n)\lqs d(x,\,x_{n_1})+d(x_{n_1},\,x_n)<\epsilon$.
\end{proof}
\end{pfbox}

\vfill\null\columnbreak

\ltitle{Complete Spaces}\v{-6}
\begin{itemize}[label=-]
\item\stitle{Definition}\\
: A metrix space $\ul X$ is \brk{a complete space} if\\
\;\;\brk{every Cauchy seq. on $\ul X$ converges in $\ul X$}.\v9

\item\stitle{Complete Subspace}\\
: For a subspace $\ul A$ of $\ul X$,\\
\;\;\brk{$\ul A$ is a complete space}\quad$\ifff$\quad\brk{$A$ is closed in $\ul X$}.

\begin{pfbox}
\begin{proof}[\color{TGreen}($\Longto$\!)]\rnc{\qedsymbol}{\Sq}Let $x\in\ol A$. Then $\exists (a_n)\subseteq A$ s.t. $(a_n)\longto x$.

Thus $(a_n)$ is a Cauchy seq. in $\ul X$ and $\ul A$.

By the assumption, $(a_n)\longto x\in A$ in $\ul A$, $\ol A\subseteq A$.
\end{proof}\v{-3}

\dotfill

\begin{proof}[\color{TGreen}($\Longfrom$)] Let $(a_n)$ be a Cauchy seq. on $\ul A$.\\
then $(a_n)$ is also a Cauchy seq. on $\ul X$,\\
hence it converges to $x\in X$ since $\ul X$ is complete.\v2

By the assumption, $x\in\ol A=A$, thus $\ul A$ is complete.
\end{proof}
\end{pfbox}\v9

\item \brk{$\ul X$ is complete} \;if and only if\\
nonempty \brk{subset seq. $(A_n)$} of $X$ s.t. $\left\{\h{-4}\begin{tabular}{l}$\ds\lim_{n\to\infty}\diam A_n=0$,\\$A_1\supseteq A_2\supseteq\cdots$\end{tabular}\right.$\\satisfies \brk{$\ds\bigcap_{n\in\NN}\ol{A_n}\neq\Emptyset$}.
\end{itemize}

\newpage

\Ltitle{3.2}{\!\begin{tabular}{l}Expansion of (Unif.) Cont. Ftn.\v{-2}\\
on Complete Space\v{-2}\end{tabular}}

\Ltitle{3.3}{Baire's Category Theorem}

\ltitle{Dense Set}\v{-6}
\begin{itemize}[label=-]
\item\stitle{Definition}\\
: For $A\subseteq X$, $A$ is \brk{dense in $\ul X$} if $\brk{\ol A=X}$.\v9

\item\stitle{Expanstion on Unif. Cont. Ftn.}\\
: For a \brk{dense subset $\ul A$} of $\ul X$ and \brk{complete space $\ul Y$},\\
\;\;every \brk{unif. cont. ftn.} $\brk f\!:\brk{\ul A}\longto\ul Y$\\
\;\;\brk{can be expanded} to a \brk{unif. cont. ftn.} $\brk g\!:\brk{\ul X}\longto\ul Y$.\v9
\end{itemize}

\ltitle{Baire's Category Theorem}\v{-6}
\begin{itemize}[label=-]
\item\stitle{Nowhere Denseness}\\
: For $A\subseteq X$, $A$ is \brk{nowhere dense in $\ul X$} if\, \brk{$\inner{\ol A}=\Emptyset$}.\\
\;\;{\Gray(ex. $\ZZ$, $\NN$: $\forall\epsilon>0$ and $\forall p\in\ZZ$, $S(p,\,\epsilon)\not\subseteq\ZZ$.\\
\qquad\;\,or every finite subset of $\RR$)}\v9

\item\stitle{1st Category Set}\\
: $A$ is a \brk{1st category set} if\\
\;\;\brk{$\exists$ nowhere dense subset seq. $(A_n)$} s.t. \brk{$A=\ds\bigcup_{n\in\NN}A_n$}.\v3
\item\stitle{2nd Category Set}\\
: $A$ is a \brk{2nd category set} if $A$ is \brk{not in 1st category}.\v9

\item\stitle{Isolated Point}\\
: $x\in X$ is an \brk{isolated point} if $\{x\}=\inner{\{x\}}$.\v9

\item For an \brk{isolated point $x$}, \brk{every $A\ni x$} is a \brk{2nd category set}.

\begin{Pfbox}
Assume $A$ is in 1st category.\\
Then $\exists$ nowhere dense subset $(A_n)$ s.t. $A=\bigcup A_n$.

If $x\in A$, $\exists n\in\NN$ s.t. $x\in A_n$, i.e., $\{x\}\subseteq A_n$.

Since $\{x\}\subseteq\ol{\{x\}}$, $\{x\}=\inner{\{x\}}\subseteq\inner{\ol{\{x\}}}=\Emptyset$: contrary.

Thus $A$ is in 2nd category.
\end{Pfbox}\v9

\item\stitle{Baire's Category Theorem}\\
: \brk{Every nonempty complete space is a 2nd category set}.\\
\;\; If \brk{$A\subseteq X$ is 1st ct.}, then \brk{$X\setminus A$ is dense in $\ul X$}.\v9

\item\stitle{Categories of Typical Sets}
\begin{itemize}[label=\Cdot]
\setlength{\leftskip}{-10pt}
\item $\brk\ZZ$, $\brk\NN$, $\brk\QQ$ are \brk{1st category sets}.\\
(Note that $\ol\ZZ=\ZZ$, $\ol\NN=\NN$, and $\inner\ZZ=\inner\NN=\Emptyset$.)\v4
\item $\RR$ is a \brk{2nd category set}.
\end{itemize}
\end{itemize}

\vfill\null\columnbreak

\Ltitle{3.4}{Topol. Complete Space}

\ltitle{Definition}\\
\quad: $\ul X$ is a \brk{topologically complete space}\\
\quad\;\;if \brk{$\ul X$ is topologically isomorphic to complete space $\ul Y$}.\v9

\ltitle{Theorems}\v{-6}
\begin{itemize}[label=-]
\item If $\ul X$ is \brk{topol. complete space}, then $\ul X$ is in \brk{2nd category}.\v6
\item For a topol. complete space $\ul X$,\\
Both \brk{closed} and \brk{open subspace} of $\ul X$ are \brk{topol. complete}.
\end{itemize}


\Hline

\Ltitle{3.5}{Completion of Metric Spaces}

\ltitle{Definition}\\
\quad: $\ul Y$ is a \brk{completion} of $\ul X$, if\v{-6}
\begin{itemize}[label=-]
\setlength{\leftskip}{4pt}
\item $\ul Y$ is a \brk{complete space};
\item \brk{$\ul X$ is a dense subset} of $\ul Y$. (i.e., $\ol X=Y$.)
\end{itemize}

\ltitle{Theorems}\v{-6}
\begin{itemize}[label=-]
\item If $\ul{Y_1}$, $\ul{Y_2}$ are completions of $\ul X$,\\
then \brk{$\exists$ an isometry} $f\!:\ul{Y_1}\longto\ul{Y_2}$, $x\longmapsto x$ ($x\in\ul X$).\v9

\item If \brk{$\ul X$ is metrically isomorphic} to a \brk{subset of a complete set},\\
then \brk{$\ul X$ has its completition}.\v9

\item \brk{Every metric space $\ul X$ has its competition $C\ul X$}.\v{15}
\end{itemize}

\newpage

\Ltitle{4.1}{Totally Bounded Metric Spaces}

\ltitle{Definitions} \;($\ul X$: metric space)\v{-6}
\begin{itemize}[label=-]
\item\stitle{$\epsilon$-dense} ($A\subseteq X$)\\
: $A$ is \brk{$\epsilon$-dense} on $\ul X$, if given $\epsilon>0$,\v{-4}
$$\brk{\dist(x,\,A)\lqs\epsilon\quad\tx{for all}\;x\in X}\tx.$$

\item\stitle{$\epsilon$-net}\\
: A \brk{$\epsilon$-net of $\ul X$} is a \brk{finite $\epsilon$-dense subset} of $\ul X$.\v9

\item\stitle{Total Boundedness}\\
: $\ul X$ is \brk{totally bounded} if\v{-4}
\begin{center}$X=\Emptyset$\quad or\quad\brk{$\exists$ an $\epsilon$-net on $\ul X$ for all $\epsilon>0$}.\end{center}

\begin{itemize}[label=\Cdot]
\setlength{\leftskip}{-10pt}
\item\stitle{Judging Theorems}\v2\\
\h{-10}\begin{tabular}{lll}$\tx:\h{-25}$&&$\ul X$ is totally bounded\v2\\
&$\ifff\h{-5}$&Every $(x_n)$ on $\ul X$ has \brk{Cauchy subseq.}\v2\\
&$\ifff\h{-5}$&Every $(x_n)$ on $\ul X$ has \brk{an acc. pt. on $C\ul X$.}
\end{tabular}\v9
\end{itemize}
\end{itemize}

\ltitle{Theorems}\v{-6}
\begin{itemize}[label=-]
\item\stitle{Total Boundedness of Range}\\
: If $f\!:\ul X\longto\ul Y$ is \brk{surjective unif. cont. ftn.}\\
\;\;where \brk{$\ul X$ is totally bounded},\v2\\
\;\;then \brk{$\ul Y$ is also totally bounded}.\v9

\item \brk{Every totally bounded metric space is bounded}.\v9

\item For \brk{unif. cont.} $f\!:\ul X\longto\ul Y$, if \brk{$\ul X$ is totally bounded},\\
then \brk{$\ul{f(X)}$ is [totally] bounded in $\ul Y$}.\v9

\item\stitle{Total Boundedness of Subspaces} \;($\ul A$: subspace of $\ul X$)
\begin{itemize}[label=\Cdot]
\setlength{\leftskip}{-10pt}
\item If \brk{$\ul X$ is totally bounded}, then \brk{so is $\ul A$}.\v4
\item If \brk{$A$ is dense in $\ul X$} and \brk{$\ul A$ is totally bounded},\\
then \brk{$\ul X$ is also totally bounded}.\v9
\end{itemize}

\item\stitle{Total Boundedness of Completion}\\
: \brk{$\ul X$ is totally bounded}\quad$\ifff$\quad\brk{$C\ul X$ is totally bounded}.
\end{itemize}

\Hline

\Ltitle{4.2}{Compact Metric Space}

\ltitle{Definition}\\
\quad\;\;: $\ul X$ is a \brk{compact metric space},\\
\qquad\;if every seq. $(x_n)$ in $\ul X$ \brk{has its acc. pt.}\v{12}

\h{-4}\Cdot\; \brk{$\ul X$ is compact}\quad$\ifff$\quad\brk{$\ul X$ is complete and totally bounded}.\v9

\h{-4}\Cdot\; For a \brk{subset $A$ of a compact space $\ul X$},\v{-4}
\begin{center}\brk{$\ul A$ is compact}\quad$\ifff$\quad\brk{$\ol A\subseteq X$}.\v2\end{center}

\vfill\null\columnbreak

\h{-4}\Cdot\; For subsets $A$, $B$ of compact space $\ul X$,\v{-4}
\begin{center}\brk{$A$, $B$: nbh. rel.}\quad$\ifff$\quad\brk{$\ol A\cap\ol B\neq\Emptyset$}.\end{center}

\h{-4}\Cdot\; For a compact space $\ul X$,\\
\quad\;\:if \brk{$f\!:\ul X\longto\ul X$ is cont.}, then \brk{$f$ is unif. cont.}

\begin{Pfbox}
For $A$, $B\subseteq X$ which is in nbh. rel.,\\
$\exists x\in\ol A\cap\ol B$ because $\ul X$ is complete.\v2

Since $f$ is cont., $f(x)\in\ol{f(A)}\cap\ol{f(B)}$,\\
hence $f(A)$, $f(B)$: in nbh. rel., thus $f$ is unif. cont.
\end{Pfbox}\v9

\h{-4}\Cdot\; For a compact space $\ul X$,\\
\quad\;\:if \brk{$f\!:\ul X\longto\ul Y$} is \brk{cont.} \& \brk{surj.}, then \brk{$\ul Y$: compact}.\v9

\h{-4}\Cdot\; For a compace space $\ul X$,\\
\quad\;\:\brk{any top. iso. $f\!:\ul X\longto\ul Y$ is unif. iso}.

\h{-4}\Cdot\; For a ftn. $f\!:\ul X\longto\ul Y$ with \brk{compact $\ul X$},\v3\\
\quad\begin{tabular}{ll}
&$f$ is \brk{bij. \& cont};\v1\\
$\ifff$&$f$ is a \brk{topological isomorphism};\v1\\
$\ifff$&$f$ is a \brk{uniform isomorphism}.
\end{tabular}\v{25}


\end{multicols*}
\end{document}